\documentclass[a4paper]{article}
% Español (México) con XeLaTeX: fontspec para fuentes Unicode, polyglossia
% para la división silábica y los nombres del documento en la variante mexicana.
\usepackage{fontspec}
\usepackage{polyglossia}
\setdefaultlanguage[variant=mexican]{spanish}
% Los nombres chinos van como «pinyin (original)»: xeCJK da a los caracteres
% Han su propia fuente; sin ella XeLaTeX los omite con un aviso.
% FandolSong no cubre algunos caracteres de nombres (U+73FA); WenQuanYi Zen Hei sí.
\usepackage[AutoFallBack=true]{xeCJK}
\setCJKmainfont{FandolSong-Regular.otf}
\setCJKfallbackfamilyfont{\CJKrmdefault}{WenQuanYi Zen Hei}
% polyglossia no traduce los nombres de listings.
\AtBeginDocument{\providecommand{\lstlistingname}{}\renewcommand{\lstlistingname}{Listado}%
  \providecommand{\lstlistlistingname}{}\renewcommand{\lstlistlistingname}{Índice de listados}}
\usepackage{amsmath,amssymb}
\usepackage{graphicx}
\usepackage[margin=2.35cm]{geometry}
\usepackage[most]{tcolorbox}
\usepackage{etoolbox}
\usepackage{listings}
\usepackage{xcolor}
\usepackage{hyperref}
\usepackage{booktabs,longtable,tabularx,array}
\usepackage{subcaption}
\usepackage{float}
\usepackage{enumitem}
\usepackage{microtype}
\usepackage{fancyhdr}
\usepackage{needspace}

\definecolor{kbBack}{HTML}{F2F6FA}
\definecolor{kbFrame}{HTML}{9DB4CC}
\definecolor{kbTitle}{HTML}{52708F}
\definecolor{ibBack}{HTML}{FBF5E7}
\definecolor{ibFrame}{HTML}{C9A961}
\definecolor{ibTitle}{HTML}{7D5F1E}
\definecolor{wbBack}{HTML}{FCF2F0}
\definecolor{wbFrame}{HTML}{D6A096}
\definecolor{wbTitle}{HTML}{9E4F43}
\definecolor{dbBack}{HTML}{F4F7F3}
\definecolor{dbFrame}{HTML}{AEC2AC}
\definecolor{dbTitle}{HTML}{5F7F64}

\tcbset{softbox/.style={
  enhanced,breakable,boxrule=0.8pt,arc=2pt,left=3mm,right=3mm,
  top=2mm,bottom=2mm,coltitle=white,fonttitle=\bfseries,
  toptitle=0.6mm,bottomtitle=0.6mm
}}
\newtcolorbox{knowledgebox}[1]{softbox,colback=kbBack,colframe=kbFrame,colbacktitle=kbTitle,title={#1}}
\newtcolorbox{importantbox}[1]{softbox,colback=ibBack,colframe=ibFrame,colbacktitle=ibTitle,title={#1}}
\newtcolorbox{warningbox}[1]{softbox,colback=wbBack,colframe=wbFrame,colbacktitle=wbTitle,title={#1}}
\newtcolorbox{dialoguebox}[1]{softbox,colback=dbBack,colframe=dbFrame,colbacktitle=dbTitle,title={#1}}

\lstset{
  language=Python,basicstyle=\ttfamily\small,
  keywordstyle=\color{kbTitle}\bfseries,stringstyle=\color{wbTitle},
  commentstyle=\color{dbTitle},breaklines=true,frame=single,numbers=left,
  numberstyle=\tiny\color{gray},captionpos=b,columns=fullflexible,
  keepspaces=true,showstringspaces=false,extendedchars=true
}
\BeforeBeginEnvironment{lstlisting}{\Needspace{12\baselineskip}}

\hypersetup{colorlinks=true,linkcolor=kbTitle,urlcolor=kbTitle,citecolor=wbTitle}
\setlist{itemsep=0.25em,topsep=0.4em}
\setlength{\parindent}{2em}
\setlength{\parskip}{0.25em}
\newcolumntype{L}[1]{>{\raggedright\arraybackslash}p{#1}}

\providecommand{\notetitle}{Notas del curso: [tema del curso]}
\providecommand{\noteauthors}{Elaboradas a partir de los materiales públicos de Stanford CS25}
\providecommand{\notedate}{Versión reescrita en 2026}
\providecommand{\videochannel}{Stanford Online}
\providecommand{\videopublishdate}{2022}
\providecommand{\videoduration}{[duración del video]}
\providecommand{\videourl}{}
\providecommand{\videocoverpath}{cover.jpg}
\providecommand{\coursesite}{https://web.stanford.edu/class/cs25/}
\providecommand{\repourl}{https://github.com/hqhq1025/ai-course-notes}
\providecommand{\skillversion}{wdkns/wdkns-skills@39f1a04}

\newcommand{\figsource}[1]{\par\vspace{-0.3em}{\footnotesize\color{gray}\emph{Fuente: #1}}\par}
\newcommand{\teachervoice}[1]{\begin{knowledgebox}{Nota de clase}#1\end{knowledgebox}}
\newcommand{\term}[1]{\textbf{#1}}

\pagestyle{fancy}
\fancyhf{}
\fancyfoot[C]{\thepage}
\fancyfoot[R]{\footnotesize\color{gray}\href{\repourl}{AI Course Notes}}
\renewcommand{\headrulewidth}{0pt}
\renewcommand{\footrulewidth}{0pt}

\newcommand{\makecscover}{
\begin{titlepage}
\centering
\vspace*{0.7cm}
{\Large Stanford CS25 · Transformers United\par}
\vspace{0.8cm}
{\huge\bfseries \notetitle\par}
\vspace{0.6cm}
{\large \noteauthors\par}
\vspace{0.25cm}
{\large \notedate\par}
\vspace{0.9cm}
\ifdefempty{\videocoverpath}{}{
  \includegraphics[width=0.86\textwidth,height=0.43\textheight,keepaspectratio]{\videocoverpath}\par
}
\vfill
\begin{tcolorbox}[width=0.92\textwidth,colback=black!2!white,colframe=black!45,boxrule=0.7pt,arc=2pt]
\textbf{Autor o canal del video}: \videochannel\par
\textbf{Fecha de publicación}: \videopublishdate\par
\textbf{Duración del video}: \videoduration\par
\textbf{Sitio del curso}: \href{\coursesite}{\nolinkurl{\coursesite}}\par
\ifdefempty{\videourl}{}{\textbf{Enlace del video}: \href{\videourl}{\nolinkurl{\videourl}}\par}
\textbf{Normas de elaboración}: \skillversion\ + las normas source-first / teacher-voice / visual-QA de este repositorio
\end{tcolorbox}
\end{titlepage}
\tableofcontents
\newpage
}

\usepackage{multicol}

\renewcommand{\notetitle}{Lecture 26: Overview of Transformers}
\renewcommand{\noteauthors}{Elaborado a partir de la clase de Div Garg, Steven Feng, Emily Bunnapradist y Seonghee Lee, y del deck oficial de 114 páginas}
\renewcommand{\notedate}{CS25 V4 · versión reescrita source-first de 2026}
\renewcommand{\videochannel}{Stanford Online}
\renewcommand{\videopublishdate}{Clase: 2024-04-04; publicación: 2024-04-23}
\renewcommand{\videoduration}{1:17:28}
\renewcommand{\videourl}{https://www.youtube.com/watch?v=fKMB5UlVY1E}
\renewcommand{\videocoverpath}{cover.jpg}

\newcommand{\lecturefigure}[3]{%
\begin{figure}[H]
  \centering
  \includegraphics[width=0.97\textwidth,height=0.50\textheight,keepaspectratio]{slides-images/#1}
  \caption{#2}
  \figsource{#3}
\end{figure}
}

\let\standardtableofcontents\tableofcontents
\renewcommand{\tableofcontents}{%
  \begingroup
  \small
  \setlength{\parskip}{0pt}
  \standardtableofcontents
  \endgroup
}

\begin{document}
\makecscover

\section{Ubicación del curso e historia del NLP: el Transformer no surgió de la nada}

Esta clase inaugural de la V4 cumple dos tareas. Primero, establece un vocabulario común para las conferencias de los invitados que siguen: attention, large language model, alignment, reasoning y agent. Segundo, presenta una línea de investigación que todavía no se cierra: las capacidades de los modelos se expanden con rapidez, pero la eficiencia en el uso de datos, el aprendizaje continuo, la memoria, la confiabilidad y los límites de seguridad siguen mostrando carencias evidentes. Al leer, conviene separar los ejemplos de productos de la clase de 2024 de los mecanismos más estables: los primeros son evidencia de un momento determinado; solo los segundos son conocimiento transferible.

\lecturefigure{slide-001.jpg}{Portada de CS25 Transformers United V4}{Slide deck oficial de CS25 V4, página 1}

\begin{knowledgebox}{Lectura de la figura: el Overview no es «repasar nombres rápidamente»}
La portada enumera a los cuatro organizadores del curso, y el contenido posterior, en efecto, lo desarrollan distintos ponentes por segmentos. Abarca historia, mecanismos, aplicaciones, debilidades, razonamiento y agents, por lo que funciona más bien como un mapa del curso. Las notas complementan las fórmulas y los límites de los sistemas, pero no reintroducen dashboards, incident runbooks ni procesos de cumplimiento normativo que no aparecieron en la clase.
\end{knowledgebox}

\teachervoice{La apertura explica que el objetivo de CS25 no es repetir los libros de texto, sino que quienes están haciendo investigación expliquen directamente los resultados nuevos, de modo que los estudiantes puedan llevar esas ideas a su propia investigación o establecer colaboraciones. Este enfoque determina que la clase presente lado a lado los consensos, los problemas abiertos y los juicios personales de los ponentes.}

\subsection{Qué espera el curso que se lleven los estudiantes}

Una visión general técnica tiende a degenerar en una lista de nombres de modelos. Los objetivos de aprendizaje que plantea el curso son más concretos: entender cómo funciona el Transformer, entender por qué logró extenderse más allá del NLP, identificar las líneas de investigación de frontera y seguir cuestionando las debilidades que persisten. Este marco exige explicar tanto las condiciones del éxito como los límites del fracaso. Al leer la página de objetivos que sigue, observe primero que los cuatro verbos apuntan, respectivamente, a mechanism, application, research y limitation, en lugar de tomarla como una página promocional del curso.

\lecturefigure{slide-011.jpg}{Los cuatro tipos de objetivos de aprendizaje del curso}{Slide deck oficial de CS25 V4, página 11}

\begin{importantbox}{Los cuatro ejes de esta clase}
\begin{enumerate}
  \item \textbf{Mechanism}: cómo transmiten información attention, self-attention, multi-head y cross-attention;
  \item \textbf{Scale and training}: cómo moldean las capacidades los parámetros, los datos, el alignment y la preference optimization;
  \item \textbf{Adaptation}: cómo la memoria, la recuperación de información, el aprendizaje continuo, el model editing y el reasoning scaffold cubren las carencias del modelo;
  \item \textbf{System}: cómo agent, tool, memory, multi-agent communication y safety convierten el modelo en un sistema ejecutable.
\end{enumerate}
\end{importantbox}

\subsection{Del rule system a la línea de tiempo de attention}

El curso presenta primero una línea de tiempo, no para discutir cuál fue el «primer artículo» de cierta técnica, sino para mostrar qué cuello de botella corregía cada generación de métodos. El NLP rule-based escribe el lenguaje como reglas; el word embedding asigna a cada palabra discreta un vector continuo; las RNN/LSTM procesan secuencias; attention permite que el modelo asigne pesos distintos a distintas posiciones; el Transformer convierte attention en la estructura de cómputo principal.

\lecturefigure{slide-013.jpg}{Línea de tiempo desde los sistemas de reglas hasta el Transformer y los modelos multimodales}{Slide deck oficial de CS25 V4, página 13}

\begin{knowledgebox}{Términos clave: los cinco escalones de la línea de tiempo}
\begin{tabularx}{\textwidth}{L{0.19\textwidth} X}
\toprule
Etapa & Problema principal que resuelve \\
\midrule
Rule-based system & Codifica la sintaxis y la semántica con reglas manuales; es transparente, pero su cobertura es frágil y su costo de transferencia es alto. \\
Word embedding & Expresa la similitud con vectores continuos, lo que permite al modelo compartir la estructura estadística entre palabras. \\
RNN/LSTM & Procesa secuencias de longitud variable con un recurrent state, pero el entrenamiento es secuencial y la información a larga distancia tiende a atenuarse. \\
Attention & Pondera las posiciones de la entrada según su contenido, sin comprimir toda la historia en un único estado. \\
Transformer & Usa attention y redes de propagación hacia adelante como columna vertebral, lo que amplía el margen para el entrenamiento en paralelo y el escalamiento. \\
\bottomrule
\end{tabularx}
\end{knowledgebox}

\lecturefigure{slide-014.jpg}{Ambigüedad, mantenimiento de reglas y generalización: los problemas del NLP temprano}{Slide deck oficial de CS25 V4, página 14}

\begin{warningbox}{La página histórica no respalda que «los métodos antiguos sean completamente inútiles»}
Las reglas, los diccionarios, la recuperación de información, las máquinas de estados finitos y las características tradicionales siguen siendo adecuados para tareas muy restringidas. La ventaja del Transformer está en el aprendizaje de representaciones de extremo a extremo y en el escalamiento, no en invalidar todos los métodos estructurados. Los sistemas de ingeniería suelen usar el learned model junto con reglas, bases de datos y verificadores.
\end{warningbox}

\lecturefigure{slide-015.jpg}{Los métodos de NLP pasan de las reglas escritas a mano a las representaciones aprendidas}{Slide deck oficial de CS25 V4, página 15}

\lecturefigure{slide-016.jpg}{ELIZA muestra las capacidades y los límites de los sistemas de diálogo rule-based}{Slide deck oficial de CS25 V4, página 16}

\begin{knowledgebox}{Lectura de la figura: por qué ELIZA parece «estar conversando»}
ELIZA usa coincidencia de patrones y reescritura con plantillas para insertar fragmentos de la oración del usuario en respuestas predefinidas. Esto muestra que la fluidez superficial no equivale a comprensión interna: si la forma de la interacción es lo bastante adecuada, incluso un conjunto limitado de reglas puede producir una sensación de coherencia. Hoy, al evaluar un LLM, tampoco se puede equiparar la naturalidad del lenguaje con la exactitud factual, la corrección del razonamiento o la confiabilidad de las acciones.
\end{knowledgebox}

\lecturefigure{slide-017.jpg}{Reglas lingüísticas, semantic parsing y las primeras representaciones simbólicas}{Slide deck oficial de CS25 V4, página 17}

\begin{importantbox}{El conocimiento estructural sigue teniendo valor}
El semantic parsing asigna al lenguaje natural una forma lógica o una representación ejecutable. Los LLM modernos pueden reducir las reglas escritas a mano, pero no eliminan la necesidad de estructura: el schema de tool calling, SQL, los programas, los pasos de una demostración y los agent protocols reintroducen, todos ellos, estructuras verificables.
\end{importantbox}

\lecturefigure{slide-018.jpg}{El word embedding asigna a las palabras discretas vectores en un espacio continuo}{Slide deck oficial de CS25 V4, página 18}

\begin{importantbox}{Definición mínima de embedding}
Sea $e_w\in\mathbb{R}^{d}$ el vector de la entrada $w$ del vocabulario. El embedding (incrustación) asigna a cada símbolo discreto un punto de un espacio continuo de $d$ dimensiones, de modo que, bajo cierto objetivo de entrenamiento, las palabras similares obtengan representaciones cercanas. Una medida de similitud habitual es
\[
\operatorname{cos}(e_i,e_j)=\frac{e_i^\top e_j}{\lVert e_i\rVert_2\lVert e_j\rVert_2}.
\]
Aquí $e_i,e_j$ son los vectores de dos tokens u objetos; el numerador es el producto interno y el denominador elimina el efecto de la longitud de los vectores. La cercanía solo indica correlación en el espacio de representación; no equivale automáticamente a una relación factual ni a una relación causal.
\end{importantbox}

\subsection{Resumen de la sección}

La posición histórica del Transformer puede condensarse en una frase: elevó la tarea de «seleccionar la información relevante de una secuencia» a un modo de cómputo principal paralelizable, apilable y escalable. Las capacidades posteriores no provienen de una sola fórmula, sino de la evolución conjunta de architecture, data, optimization, hardware e interface.
\section{Attention y Transformer: cómo se selecciona y combina la información}

La sección anterior responde «por qué se llegó al Transformer»; esta sección responde «qué hace exactamente». La clase parte de seq2seq, pasa a Q/K/V y luego distingue entre self-attention, multi-head attention y cross-attention. Entender estas diferencias permite ubicar más adelante la interfaz de los LLM decoder-only, de encoder--decoder, de retrieval y de la memoria de un agent.

\subsection{El cuello de botella de seq2seq: un solo state carga con demasiado}

Los primeros modelos sequence-to-sequence (de secuencia a secuencia) usaban un encoder para comprimir la entrada en un estado oculto, a partir del cual un decoder generaba la salida paso a paso. Al aumentar la longitud, un único vector debe conservar al mismo tiempo el significado de las palabras, el orden y las dependencias lejanas, y el cuello de botella de información se vuelve cada vez más evidente. Las dos diapositivas siguientes muestran primero el recurrent path y después la vía de lectura directa que agrega attention; lo que conviene comparar es desde dónde fluye la información hacia el decoder.

\lecturefigure{slide-019.jpg}{Un RNN seq2seq transmite la información de la secuencia mediante el recurrent state}{slide deck oficial de CS25 V4, página 19}

\begin{importantbox}{Recurrent update}
El RNN más simple puede escribirse como
\[
h_t=\phi(W_xx_t+W_hh_{t-1}+b),\qquad y_t=W_oh_t.
\]
$x_t$ es la $t$-ésima entrada, $h_t$ es el estado oculto, $W_x,W_h,W_o$ son matrices aprendibles y $\phi$ es una función no lineal. Cada paso depende de $h_{t-1}$, por lo que la dimensión temporal difícilmente puede paralelizarse por completo; además, en secuencias largas el gradiente debe atravesar muchos recurrent steps.
\end{importantbox}

\lecturefigure{slide-020.jpg}{Attention permite que el decoder acceda directamente a varios encoder states}{slide deck oficial de CS25 V4, página 20}

\begin{knowledgebox}{Lectura de la figura: lo que cambia attention es la vía de la información}
Sin attention, el decoder depende principalmente del encoder state final; con attention, cada posición de salida puede asignar una puntuación a todas las posiciones de entrada y calcular una suma ponderada. El modelo ya no está obligado a comprimir la oración completa en un cuello de botella fijo, sino que vuelve a leer la entrada según lo que requiera la generación en curso.
\end{knowledgebox}

\subsection{Q, K, V: consulta, índice y contenido}

En la clase se usa una analogía con un library system para explicar Query, Key y Value. El lector formula una consulta, el índice del catálogo funciona como key y el contenido del libro funciona como value. El modelo primero compara la query con cada key para medir su relevancia y después combina los value correspondientes mediante pesos. Al leer la figura hay que separar «qué libro se encontró» de «qué se extrae del libro», lo que corresponde precisamente a los dos pasos de key matching y value aggregation.

\lecturefigure{slide-021.jpg}{La library analogy explica Query, Key y Value}{slide deck oficial de CS25 V4, página 21}

\begin{importantbox}{Scaled dot-product attention}
Sean $Q\in\mathbb{R}^{n_q\times d_k}$, $K\in\mathbb{R}^{n_k\times d_k}$ y $V\in\mathbb{R}^{n_k\times d_v}$; attention se define como
\[
\operatorname{Attention}(Q,K,V)=\operatorname{softmax}\!\left(\frac{QK^\top}{\sqrt{d_k}}+M\right)V.
\]
$QK^\top$ da la puntuación de coincidencia entre cada query y cada key; $\sqrt{d_k}$ controla la escala del producto punto; $M$ es una mask opcional; softmax convierte las puntuaciones en pesos y, al final, se calcula una suma ponderada sobre $V$. Query determina «qué se busca ahora», Key determina «si soy relevante» y Value determina «qué contenido aporto cuando soy relevante».
\end{importantbox}

\begin{warningbox}{Un attention weight no es una explicación completa}
Los pesos muestran cómo cierta cabeza de cierta capa mezcla los value, pero la salida pasa además por la value projection, el residual, el MLP y las capas posteriores. No se puede afirmar, con base en un solo attention map, «por qué» el modelo tomó una decisión, y mucho menos interpretar un peso alto directamente como contribución causal.
\end{warningbox}

\lecturefigure{slide-022.jpg}{Los attention weights combinan varios value en un context vector}{slide deck oficial de CS25 V4, página 22}

\begin{lstlisting}[caption={Pseudocódigo didáctico de scaled dot-product attention}]
def attention(query, key, value, mask=None):
    scores = query @ key.transpose(-1, -2)
    scores = scores / sqrt(key.shape[-1])
    if mask is not None:
        scores = scores + mask
    weights = softmax(scores, dim=-1)
    return weights @ value, weights
\end{lstlisting}

\teachervoice{El ponente enfatiza que se trata de un «soft match»: una query no tiene que elegir un solo value, sino que puede tomar información de varias posiciones. Esta propiedad distingue a attention de la coincidencia exacta de una base de datos y explica también por qué puede expresar relaciones lingüísticas difusas y composicionales.}

\subsection{Self-attention y Multi-head attention}

Self-attention (autoatención) hace que una misma secuencia produzca a la vez $Q$, $K$ y $V$. Cada token puede leer otros tokens según su contenido; un causal decoder usa una mask para impedir la lectura de posiciones futuras. Multi-head attention (atención multicabeza), por su parte, permite que varios subespacios aprendan en paralelo distintos patrones de coincidencia. Las tres diapositivas siguientes avanzan nivel por nivel, de la conexión de una sola cabeza al block completo y a la visualización multicabeza: primero se observa el origen de la información y después la projection y la fusión.

\lecturefigure{slide-023.jpg}{Self-attention establece conexiones basadas en el contenido dentro de una misma secuencia}{slide deck oficial de CS25 V4, página 23}

\begin{knowledgebox}{Concepto previo: tres masks comunes en self-attention}
\begin{tabularx}{\textwidth}{L{0.23\textwidth} X}
\toprule
Forma & Flujo de información permitido \\
\midrule
Bidirectional & Cada posición puede leer toda la entrada; es típico del encoder. \\
Causal & La posición $t$ solo puede leer los token $\le t$; es típico del autoregressive decoder. \\
Block/sparse & Solo se permiten bloques locales, token globales o conexiones predefinidas; se usa para reducir el costo en secuencias largas. \\
\bottomrule
\end{tabularx}
\end{knowledgebox}

\lecturefigure{slide-024.jpg}{Un Transformer block apila multi-head attention y una red feedforward}{slide deck oficial de CS25 V4, página 24}

\begin{importantbox}{Multi-head attention}
La $i$-ésima head primero realiza una proyección independiente:
\[
\operatorname{head}_i=\operatorname{Attention}(QW_i^Q,KW_i^K,VW_i^V),
\]
y después todas las head se concatenan y se proyecta la salida:
\[
\operatorname{MHA}(Q,K,V)=\operatorname{Concat}(\operatorname{head}_1,\ldots,\operatorname{head}_h)W^O.
\]
$h$ es el número de head, $W_i^Q,W_i^K,W_i^V$ son las proyecciones de cada head y $W^O$ se encarga de la fusión. Tener varias cabezas no garantiza que cada head corresponda automáticamente a un significado claro, pero sí aumenta los subespacios de representación en paralelo.
\end{importantbox}

\lecturefigure{slide-025.jpg}{Distintas attention heads pueden aprender distintas relaciones posicionales}{slide deck oficial de CS25 V4, página 25}

\begin{warningbox}{«Una head se encarga de una regla gramatical» es solo una observación, no un contrato de interfaz}
Algunas head muestran patrones de posición, de correferencia o sintácticos, pero el entrenamiento no exige que una head desempeñe permanentemente un papel fijo. Tras una poda, un reentrenamiento o un cambio de modelo, los patrones pueden cambiar. El diseño de ingeniería no puede depender de que una head con cierto número conserve necesariamente cierto tipo de conocimiento.
\end{warningbox}

\subsection{Cross-attention: cómo se alinean dos secuencias}

Cross-attention (atención cruzada) toma la query y los key/value de fuentes distintas. En traducción automática, el decoder produce la query y los encoder hidden states aportan los key y value; la generación aumentada por recuperación, los modelos de visión y lenguaje y las arquitecturas encoder--decoder también usan interfaces similares.

\lecturefigure{slide-026.jpg}{Self-attention y cross-attention en traducción automática}{slide deck oficial de CS25 V4, página 26}

\begin{knowledgebox}{Glosario: no hay que confundir los tres tipos de attention}
\begin{tabularx}{\textwidth}{L{0.22\textwidth} L{0.25\textwidth} X}
\toprule
Tipo & Origen de Q/K/V & Uso principal \\
\midrule
Self-attention & La misma secuencia & Establecer dependencias dentro de la secuencia. \\
Causal self-attention & La misma secuencia, con mask del futuro & Generación autorregresiva. \\
Cross-attention & Q proviene del decoder/una modalidad; K/V provienen del encoder/otra modalidad & Generación condicionada y fusión entre modalidades y entre secuencias. \\
\bottomrule
\end{tabularx}
\end{knowledgebox}

\subsection{En qué supera el Transformer al RNN}

La última diapositiva de comparación resume sus ventajas en acceso a larga distancia, entrenamiento estable y cómputo en paralelo. Hay que añadir una limitación: la score matrix de la dense self-attention estándar es de $n\times n$, por lo que el aumento de la longitud de la secuencia implica un costo cuadrático de cómputo y de memoria. Al leer esta tabla conviene comparar a la vez el dependency path, los sequential steps y el resource cost, para no quedarse solo con «el Transformer es mejor».

\lecturefigure{slide-027.jpg}{Comparación del entrenamiento y del modelado de dependencias entre RNN y Transformer}{slide deck oficial de CS25 V4, página 27}

\begin{importantbox}{El intercambio entre complejidad y paralelismo}
Para una secuencia de longitud $n$ y dimensión oculta $d$, el cálculo principal de los score en dense attention es de aproximadamente $O(n^2d)$, y la attention matrix ocupa $O(n^2)$; en un RNN, cada paso cuesta aproximadamente $O(d^2)$ y el cómputo total crece linealmente con $n$, pero debe ejecutarse en orden temporal. El Transformer cambia más emparejamientos globales por paralelismo en el entrenamiento y rutas de información cortas.
\end{importantbox}

\teachervoice{En la clase se considera el GPU parallelism una causa importante del éxito del Transformer. Una arquitectura no solo debe tener buena capacidad expresiva, sino también ajustarse al rendimiento del hardware; este es además el punto de partida de los problemas posteriores de scale, costo y concentración de la investigación.}

\subsection{Resumen de la sección}

Attention, mediante una ponderación basada en el contenido, permite al modelo leer la información según la necesite. La autoatención procesa las relaciones dentro de la secuencia, la atención cruzada conecta fuentes distintas y el mecanismo multicabeza ofrece varios subespacios de representación. Así, el Transformer se libra del recurrent bottleneck, pero introduce el costo cuadrático en secuencias largas y problemas de interpretabilidad más complejos.
\section{Scale, Emergence y Alignment: de dónde vienen las capacidades}

Una vez entendida la attention, el curso amplía la perspectiva del mecanismo de una sola capa al entrenamiento a gran escala. Aquí hay tres líneas que no deben confundirse: cómo la escala modifica las capacidades observables, cómo la alignment modifica las preferencias de salida y cómo los modelos de producto adoptan distintas decisiones de compromiso en datos, arquitectura y receta de entrenamiento.

\subsection{Un Large Language Model es un sistema de entrenamiento ampliado}

Un Large language model (modelo de lenguaje grande, LLM) suele seguir teniendo como núcleo la next-token prediction, pero con una escala de parámetros, datos y cómputo considerablemente mayor. Su «large» no se refiere solo a los parámetros: también implica que crecen en conjunto el pipeline de datos, el entrenamiento en paralelo, la evaluación y la infraestructura de despliegue.

\lecturefigure{slide-028.jpg}{Escala, datos de entrenamiento y next-token objective de un LLM}{Diapositivas oficiales de CS25 V4, página 28}

\begin{importantbox}{Autoregressive objective}
Para una secuencia $x_{1:T}$, la pérdida de entrenamiento suele escribirse como
\[
\mathcal{L}_{\mathrm{LM}}(\theta)=-\sum_{t=1}^{T}\log p_{\theta}(x_t\mid x_{<t}).
\]
$\theta$ son los parámetros del modelo y $x_{<t}$ son los token previos. El objetivo solo exige aumentar la verosimilitud del siguiente token; seguir instrucciones, citar evidencia, rechazar operaciones peligrosas y planificar a largo plazo no son capacidades que este objetivo garantice directamente.
\end{importantbox}

\begin{warningbox}{El número de parámetros no es un estadístico suficiente de la capacidad}
Con la misma escala de parámetros, el número de token, la calidad de los datos, la architecture, el optimizer, la context length, el post-training y la contaminación de la evaluación modifican el resultado. Cuando más adelante la clase analiza Phi, BabyLM, DPO y MoE, lo que muestra es precisamente que, además de «más grande», existen muchos grados de libertad de diseño.
\end{warningbox}

\subsection{Emergent ability y las trampas de la medición}

La Emergent ability (capacidad emergente) se define como una capacidad que aparece en los modelos grandes y no en los pequeños. La figura de la clase muestra varias tareas cuyo desempeño sube de golpe cerca de un umbral de escala, pero el ponente también cita investigaciones que lo cuestionan: una metric discreta o no lineal puede hacer que una mejora gradual se vea como un salto. Las siguientes figuras deben leerse dos veces: la primera, para ver las curvas de las tareas; la segunda, para revisar la definición del eje vertical y el threshold de puntuación.

\lecturefigure{slide-029.jpg}{Varias tareas muestran cambios bruscos aparentes con la escala del modelo}{Diapositivas oficiales de CS25 V4, página 29}

\begin{knowledgebox}{Lectura de la figura: distinguir primero el eje horizontal, el eje vertical y el threshold}
El eje horizontal suele ser la escala del modelo o el cómputo de entrenamiento, y el vertical, la task metric. Si la metric solo otorga puntos cuando la respuesta es «completamente correcta», una pequeña mejora continua de la probabilidad subyacente puede, al cruzar el umbral, registrarse de pronto como 1. La figura puede demostrar una no linealidad en cierta evaluación, pero por sí sola no demuestra que dentro del modelo haya ocurrido una transición de fase discreta.
\end{knowledgebox}

\lecturefigure{slide-030.jpg}{Desempeño de tareas con Few-shot prompting según la escala}{Diapositivas oficiales de CS25 V4, página 30}

\begin{importantbox}{Un ejemplo sencillo de Metric-induced emergence}
Supongamos que la probabilidad de que el modelo dé la respuesta correcta aumenta de forma gradual con la escala $s$ como $p(s)$. Si la evaluación exige que $k$ pasos consecutivos sean todos correctos, la exactitud global es aproximadamente
\[
A(s)=p(s)^k.
\]
Aunque $p(s)$ sea gradual, $A(s)$ sube abruptamente en la región alta. Un análisis más sólido debe considerar a la vez la token likelihood, el partial credit, la calibration y las curvas bajo distintos umbrales.
\end{importantbox}

\lecturefigure{slide-031.jpg}{Varias explicaciones posibles de la emergent ability}{Diapositivas oficiales de CS25 V4, página 31}

\teachervoice{El ponente plantea explícitamente que la emergence puede deberse en parte a la metric no lineal que eligen los investigadores, y no a que la capacidad del modelo aparezca de verdad de forma repentina. Esto no niega el papel de la scale; exige incluir también los instrumentos de observación en el análisis causal.}

\lecturefigure{slide-032.jpg}{Beyond Scaling: la arquitectura, los datos y el proceso de entrenamiento también pueden cambiar las capacidades}{Diapositivas oficiales de CS25 V4, página 32}

\begin{knowledgebox}{Por qué beyond scaling merece un análisis aparte}
Una mejor architecture puede mejorar el sesgo inductivo; datos de mayor calidad aumentan el valor de aprendizaje de cada token; un optimizer y un curriculum más estables mejoran la eficiencia del entrenamiento; retrieval, tool y memory sacan de los parámetros parte de las capacidades. En conjunto determinan la densidad de capacidad y su accesibilidad.
\end{knowledgebox}

\lecturefigure{slide-033.jpg}{Preguntas planteadas a los estudiantes sobre scale, democratización y retrieval}{Diapositivas oficiales de CS25 V4, página 33}

\begin{warningbox}{Las preguntas de la clase no tienen una única respuesta}
«¿Basta con seguir ampliando los modelos?» y «¿Qué es mejor, la memoria paramétrica o el retrieval?» son ejes de diseño, no disyuntivas. La escala puede aportar representaciones generales, el retrieval proporciona evidencia actualizable, los modelos pequeños reducen costos y las estructuras especializadas aumentan la eficiencia. Los sistemas reales suelen combinar estas vías.
\end{warningbox}

\subsection{RLHF y DPO: de la likelihood a la preference}

El preentrenamiento hace que el modelo se ajuste a la distribución del corpus; el post-training acerca las salidas a las preferencias humanas y a las especificaciones de la tarea. El Reinforcement Learning from Human Feedback (aprendizaje por refuerzo a partir de retroalimentación humana, RLHF) suele entrenar un reward model y luego optimizar la policy con aprendizaje por refuerzo; la Direct Preference Optimization (optimización directa de preferencias, DPO) construye directamente, a partir de pares de preferencias, un objetivo de tipo clasificación.

\lecturefigure{slide-035.jpg}{Preference collection, reward model y policy optimization en RLHF}{Diapositivas oficiales de CS25 V4, página 35}

\begin{importantbox}{Forma básica de los Preference data}
Para un prompt $x$, una persona o una regla compara dos respuestas $y_w$ y $y_l$, donde $y_w$ es la preferred response. El Reward model suele aprender
\[
P(y_w\succ y_l\mid x)=\sigma\!\left(r_\phi(x,y_w)-r_\phi(x,y_l)\right).
\]
$r_\phi$ es el reward score y $\sigma$ es la logistic function. Los datos solo expresan preferencias relativas, no equivalen a la verdad objetiva; la composición de los anotadores, la rubric y la estrategia de muestreo se trasladan al comportamiento del modelo.
\end{importantbox}

\lecturefigure{slide-036.jpg}{DPO optimiza directamente el preference pair con una reference policy}{Diapositivas oficiales de CS25 V4, página 36}

\begin{importantbox}{Intuición del DPO objective}
Una forma común de escribirlo es
\[
\mathcal{L}_{\mathrm{DPO}}=-\mathbb{E}\log\sigma\!\left(\beta\left[
\log\frac{\pi_\theta(y_w\mid x)}{\pi_{\mathrm{ref}}(y_w\mid x)}-
\log\frac{\pi_\theta(y_l\mid x)}{\pi_{\mathrm{ref}}(y_l\mid x)}
\right]\right).
\]
$\pi_\theta$ es la policy por entrenar, $\pi_{\mathrm{ref}}$ es el reference model y $\beta$ controla qué tanto se permite alejarse de la reference. DPO simplifica el proceso de entrenamiento, pero el preference bias, la coverage y la over-optimization siguen presentes.
\end{importantbox}

\teachervoice{La clase no describió DPO como «RLHF ya quedó obsoleto»; subrayó que RLHF depende de retroalimentación de alta calidad, del reward y de la policy optimization, y que DPO ofrece una vía de preference-learning más directa.}

\subsection{ChatGPT, GPT-4 y Gemini: una instantánea de la clase de 2024}

Las siguientes diapositivas de modelos pertenecen al contexto del curso de 2024. Sirven para ilustrar direcciones como el post-training, la multimodalidad, el context, los benchmark y MoE, y no deben interpretarse como una clasificación de los productos más recientes de 2026. El orden de lectura debe partir de las etapas de entrenamiento y de la architecture, y después considerar los benchmark de entonces, en lugar de presuponer conclusiones a partir del nombre de la marca.

\lecturefigure{slide-037.jpg}{Panorama de ChatGPT y la serie GPT en la clase}{Diapositivas oficiales de CS25 V4, página 37}

\begin{knowledgebox}{Lectura de la figura: detrás del nombre del producto hay que separar las etapas de entrenamiento}
Un chat model suele pasar por un preentrenamiento general, un instruction tuning supervisado, preference optimization, datos de seguridad y políticas de despliegue. Que la interfaz sea parecida no significa que los datos de entrenamiento, la estructura del modelo o los límites de su comportamiento sean iguales; cuando la información pública es insuficiente, no debe deducirse todo el mecanismo interno a partir de la experiencia de uso del producto.
\end{knowledgebox}

\lecturefigure{slide-038.jpg}{Resumen de la clase sobre las capacidades, la evaluación y la multimodalidad de GPT-4}{Diapositivas oficiales de CS25 V4, página 38}

\begin{warningbox}{Un Benchmark snapshot no es una clasificación permanente}
Las puntuaciones dependen de la versión, el prompt, el conjunto de evaluación, el grado de contaminación y el judge. Los resultados presentados en clase deben conservar su fecha; no se puede usar un benchmark antiguo para afirmar que hoy cierto modelo es absolutamente superior. Lo más importante es entender las dimensiones de la evaluación: conocimiento, reasoning, multimodality, safety y cost no quedan cubiertos por una sola puntuación total.
\end{warningbox}

\lecturefigure{slide-039.jpg}{La serie Gemini y la comparación de modelos de 2024}{Diapositivas oficiales de CS25 V4, página 39}

\lecturefigure{slide-040.jpg}{La diapositiva de Gemini usa MoE para explicar los routed experts}{Diapositivas oficiales de CS25 V4, página 40}

\begin{importantbox}{Mecanismo mínimo de Mixture of Experts}
La Mixture of Experts (mezcla de expertos, MoE) usa un router para elegir unos pocos expert para cada token:
\[
p_i(x)=\operatorname{softmax}(W_rx)_i,\qquad
h'=\sum_{i\in\operatorname{TopK}(p(x))}p_i(x)E_i(x).
\]
$W_r$ son los parámetros del router, $p_i$ es el peso del expert y $E_i$ es el $i$-ésimo expert. La Sparse activation permite que, al aumentar el número total de parámetros, cada token solo se calcule con unos pocos expert; el costo es que el load balancing, la comunicación y la gestión de la capacity se vuelven más complejos.
\end{importantbox}

\begin{warningbox}{Un «Expert» no equivale a un experto humano identificable}
Un expert puede mostrar preferencia por ciertos temas o funciones, pero el entrenamiento normalmente no garantiza que se encargue solo de imágenes, matemáticas o un idioma específico. Representar MoE directamente como una división estable de oficios es solo una intuición y no sustituye el análisis empírico del routing y de la activation.
\end{warningbox}

\subsection{La situación en 2024 y las carencias por resolver}

El curso resume el estado de entonces en dos diapositivas: los LLM, la alignment, la multimodalidad y los diffusion transformer avanzaban con rapidez; los siguientes pasos incluían agent más generales, eficiencia, modelos del mundo, aprendizaje continuo y control de seguridad. En esta sección, este grupo de figuras sirve de transición entre capítulos: primero distingue las capacidades que ya se convirtieron en productos y después señala las carencias que siguen siendo objetivos de investigación, con lo que se arma una lista de problemas para las aplicaciones en otros dominios y las limitaciones de los sistemas que se tratan más adelante.

\lecturefigure{slide-041.jpg}{Where we are: síntesis de la clase de 2024 sobre el ecosistema de los LLM}{Diapositivas oficiales de CS25 V4, página 41}

\lecturefigure{slide-042.jpg}{The Future: aplicaciones más amplias, agent y eficiencia de los modelos}{Diapositivas oficiales de CS25 V4, página 42}

\begin{knowledgebox}{Lectura de la figura: la diapositiva Future mezcla tendencias técnicas y visiones de investigación}
Algunas direcciones ya cuentan con prototipos de sistemas, como la multimodal generation, el tool use y los smaller models; otras siguen siendo objetivos, como agent autónomos estables, aprendizaje continuo de bajo costo y una eficiencia de datos más cercana a la humana. Estas notas separan lo «ya demostrado» del «despliegue confiable».
\end{knowledgebox}

\lecturefigure{slide-043.jpg}{Missing ingredients: los problemas pendientes hacia sistemas más generales}{Diapositivas oficiales de CS25 V4, página 43}

\begin{importantbox}{Seis carencias tras la ampliación de capacidades}
\begin{enumerate}
  \item menor costo de entrenamiento e inferencia;
  \item representaciones más robustas del mundo real y de las relaciones causales;
  \item memoria de largo plazo actualizable y aprendizaje continuo;
  \item mecanismos internos que puedan inspeccionarse y editarse;
  \item error correction en trayectorias largas;
  \item permisos de usuario, security y value alignment.
\end{enumerate}
\end{importantbox}

\subsection{Resumen de la sección}

La Scale puede ampliar las representaciones y la cobertura de tareas, pero la emergent curve depende de la metric; la preference optimization cambia el comportamiento, pero hereda los sesgos de los datos y del objetivo; MoE, retrieval, data quality y specialized training ofrecen palancas más allá de la scale. Lo más valioso de las diapositivas de modelos no son las clasificaciones antiguas, sino la descomposición de las capacidades en etapas de entrenamiento y componentes del sistema.
\section{Aplicaciones de Transformer en distintos dominios: el mismo patrón de cómputo con distintas interfaces de datos}

En siete diapositivas, la clase muestra cómo Transformer se extendió del texto a la voz, la música, la visión, la robótica, los juegos y la biología. La conclusión que sí se puede trasladar no es que «todas las tareas quedaron unificadas», sino que en muchos dominios los objetos se pueden representar como token o set, y después se usa attention para modelar interacciones a larga distancia. La representación de los datos, las restricciones sobre la salida y los criterios de evaluación siguen siendo muy propios de cada dominio.

\subsection{Texto y lenguaje: la token interface más madura}

El texto es, por naturaleza, una secuencia discreta, por lo que es el que se conecta con más facilidad a un autoregressive objective. Traducción, resumen, preguntas y respuestas, código y diálogo parecen tareas distintas, pero con frecuencia todas pueden reformularse como token generation condicionada. La siguiente diapositiva de aplicaciones debe desglosarse según la entrada, la salida y el criterio de verificación; no se puede suponer que las tareas tienen la misma estructura solo porque todas producen texto.

\lecturefigure{slide-045.jpg}{Aplicaciones de Transformer en la generación y comprensión de texto y en el diálogo}{Presentación oficial de CS25 V4, diapositiva 45}

\begin{knowledgebox}{Lectura de la figura: una interfaz unificada no es una tarea unificada}
Escribir una tarea como text-to-text permite compartir el modelo y el código de entrenamiento, pero la loss, el decoding, la exactitud factual y los objetivos de evaluación siguen siendo distintos. La traducción se ocupa de conservar el significado; el resumen, de la cobertura y la compresión; el código, de que sea ejecutable; y el diálogo, además, debe manejar el contexto, la personalidad y las normas de seguridad.
\end{knowledgebox}

\subsection{Audio: la forma de onda debe convertirse primero en una representación modelable}

El audio es una señal continua de alta frecuencia. Por lo general, el sistema primero convierte la forma de onda en un spectrogram, en codec token o en learned feature, y después usa Transformer para modelar las relaciones temporales; la salida puede ser texto, token acústicos o eventos musicales. Al leer la figura, primero identifique la signal representation y después determine si el modelo está reconociendo, generando o realizando una alineación entre modalidades.

\lecturefigure{slide-046.jpg}{Panorama de aplicaciones de Speech Transformer y Music Transformer}{Presentación oficial de CS25 V4, diapositiva 46}

\teachervoice{Este grupo de application slide solo sirve como recorrido de orientación; el expositor deja los mecanismos a fondo para los invitados posteriores. Por eso, estas notas explican la interfaz común y las diferencias entre dominios, y no amplían el nombre de cada modelo en conclusiones sobre productos que la clase no respalda.}

\begin{importantbox}{Los tres pasos comunes del modelado de modalidades}
\begin{enumerate}
  \item \textbf{Tokenizer/encoder}: convierte la señal original en una representación de longitud finita;
  \item \textbf{Backbone}: usa attention para modelar las relaciones condicionales entre representaciones;
  \item \textbf{Decoder/renderizado}: convierte el hidden state de vuelta en texto, imagen, sonido o acción.
\end{enumerate}
Lo que se unifica es el cómputo intermedio; lo más difícil suele estar en la representación y la evaluación de los dos extremos.
\end{importantbox}

\subsection{Vision: patch, frame y alineación entre modalidades}

Vision Transformer (Transformer visual, ViT) divide la imagen en patch token; en el video, además, hay que codificar la dimensión temporal. Las tareas de análisis producen label, box o texto, mientras que las tareas de generación necesitan decodificar latent/token en una secuencia de píxeles. Las dos diapositivas siguientes corresponden a perception y a generation; lo que conviene comparar es la tokenización, el decoder y el objetivo de evaluación.

\lecturefigure{slide-047.jpg}{Vision Transformer analiza imágenes y video con patch token}{Presentación oficial de CS25 V4, diapositiva 47}

\begin{importantbox}{Patch embedding}
Si la imagen mide $H\times W$ y el patch mide $P\times P$, el número de token es aproximadamente
\[
N=\frac{HW}{P^2}.
\]
Un patch más pequeño conserva más detalle local, pero hace que el costo $N^2$ de attention crezca con rapidez. Por eso, los sistemas de visión suelen usar estructuras jerárquicas, attention local, downsampling o un latent bottleneck.
\end{importantbox}

\lecturefigure{slide-048.jpg}{Diffusion y Transformer para la generación de imágenes y video}{Presentación oficial de CS25 V4, diapositiva 48}

\begin{warningbox}{Que la imagen generada se vea real no significa que el modelo del mundo sea correcto}
Los modelos de imagen y video pueden producir texturas de alta fidelidad y, aun así, violar la permanencia de los objetos, la geometría, la causalidad o la coherencia a lo largo del tiempo. La metric de calidad visual y la corrección física no son el mismo objetivo; no se puede sustituir la verificación de la estructura y la dinámica por «se parece».
\end{warningbox}

\subsection{Robotics, juegos y biología: la salida debe someterse a restricciones externas}

La robótica y los juegos asignan la salida a una action; los modelos biológicos asignan secuencias o estructuras tridimensionales a funciones y predicciones. En estos casos, un error ya no consiste solo en generar una oración mala, sino que puede cambiar una trayectoria física, una estrategia o una decisión experimental. Las tres diapositivas siguientes deben leerse a partir de la retroalimentación del entorno y del costo de verificación, y no como si los modelos de distintos dominios fueran solo un mosaico de aplicaciones.

\lecturefigure{slide-049.jpg}{Transformer en robótica, simulación y physical task}{Presentación oficial de CS25 V4, diapositiva 49}

\begin{knowledgebox}{Términos clave: policy y world feedback}
Una policy (política) asigna una observation a una action. Un modelo de lenguaje suele generar los token de una sola vez, mientras que la policy de un robot recibe continuamente la retroalimentación de sus sensores y actualiza sus acciones. La latencia, la acumulación de errores, las restricciones de colisión y los mecanismos de recuperación dentro del lazo cerrado hacen que un embodied system sea más difícil de validar que la generación fuera de línea.
\end{knowledgebox}

\lecturefigure{slide-050.jpg}{Aplicaciones de Transformer en go, estrategia en tiempo real y entornos de juego}{Presentación oficial de CS25 V4, diapositiva 50}

\lecturefigure{slide-051.jpg}{Modelos de lenguaje médicos, estructura de proteínas y aplicaciones a secuencias biológicas}{Presentación oficial de CS25 V4, diapositiva 51}

\begin{warningbox}{Una lista de nombres por dominio debe remitirse a la tarea y a la evidencia}
Med-PaLM, AlphaFold, robotics transformer y game agent resuelven objetos distintos, con señales de supervisión y requisitos de seguridad completamente diferentes. Más que memorizar los nombres, conviene recordar tres cosas: cómo se representa la entrada, cómo se verifica la salida y quién asume las consecuencias de un error.
\end{warningbox}

\subsection{Resumen de la sección}

La expansión de Transformer a otros dominios proviene de la capacidad general de attention para modelar relaciones en conjuntos y secuencias, pero cada dominio sigue necesitando su propio tokenizer, decoder, constraint y evaluation. Un backbone unificado reduce el costo de diseñar modelos, pero no elimina la necesidad del conocimiento del dominio ni de la verificación en lazo cerrado.
\section{Limitaciones y adaptación de los LLM: datos, memoria, actualización e interpretabilidad}

Cuanto más amplias son las aplicaciones, más evidentes se vuelven las carencias del modelo. Esta sección sigue el orden de la clase y aborda cinco grupos de problemas: costo de entrenamiento y eficiencia de datos; memoria y personalización; datos sintéticos y memorization; aprendizaje continuo y model editing; self-reflection y hallucination. Todos responden a la misma pregunta: una vez congelados los parámetros, ¿cómo sigue aprendiendo, actualizándose y corrigiéndose el sistema?

\subsection{Datos, cómputo y concentración de la investigación}

Entrenar modelos grandes requiere enormes cantidades de tokens, accelerator time, energía e infraestructura de ingeniería. La clase vincula este punto con la democratización de la investigación: cuando solo unas cuantas instituciones pueden costear experimentos reproducibles, se restringen la comparación de métodos, la auditoría abierta y la participación de los estudiantes. Esta subsección revisa primero la diapositiva de costos y después pasa a BabyLM y a los modelos pequeños, para responder si «la mejora de capacidades solo puede comprarse con más cómputo».

\lecturefigure{slide-053.jpg}{Demanda de datos, cómputo, financiamiento y energía en el entrenamiento de LLM}{slide deck oficial de CS25 V4, página 53}

\begin{importantbox}{Escala aproximada del cómputo de entrenamiento}
Para el preentrenamiento de un Dense Transformer suele usarse la aproximación
\[
C\approx 6ND,
\]
donde $N$ es el número de parámetros, $D$ es el número de tokens de entrenamiento y $C$ es el orden de magnitud de las operaciones de punto flotante. La constante varía según la architecture y la implementación, pero la aproximación revela dos hechos: tanto los parámetros como los datos elevan linealmente el cómputo de entrenamiento; y tokens de mayor calidad, activación dispersa y kernels eficientes pueden cambiar el costo real.
\end{importantbox}

\begin{warningbox}{El costo no es solo una factura de GPU}
La recolección y el filtrado de datos, los experimentos fallidos, los checkpoints, la comunicación de red, el servicio de inferencia, la evaluación y la retroalimentación humana consumen recursos. Reportar solo los FLOPs del entrenamiento final subestima la cadena completa de investigación y desarrollo, y tampoco indica el costo por unidad de capacidad.
\end{warningbox}

\subsection{BabyLM: qué pregunta plantea la eficiencia de datos humana}

La diapositiva de BabyLM compara a los niños con los LLM: los niños están expuestos a mucho menos lenguaje, pero se apoyan en la percepción, la interacción, la retroalimentación social y la experiencia corporizada para formar conceptos sólidos. Esta analogía no demuestra que las máquinas deban replicar el aprendizaje infantil, pero sí expone el problema de eficiencia de datos que tiene el escalamiento basado solo en texto. Al leer las dos diapositivas en paralelo, conviene atender a la supervision density y a la interaction, y no solo comparar el número de palabras.

\lecturefigure{slide-054.jpg}{Diferencias en el entorno de datos entre el aprendizaje infantil y el entrenamiento de LLM}{slide deck oficial de CS25 V4, página 54}

\lecturefigure{slide-055.jpg}{Volumen de datos, proceso de desarrollo y experiencia multimodal en la discusión sobre BabyLM}{slide deck oficial de CS25 V4, página 55}

\begin{knowledgebox}{Lectura de la figura: los niños no son «modelos de lenguaje con pocos datos»}
Los niños cuentan con visión, acción, atención conjunta, retroalimentación emocional, intervención causal y un entorno persistente. Comparar directamente el número de palabras con los tokens de un LLM pasa por alto la supervision density y la interaction. El valor de BabyLM está en obligar a los investigadores a explorar mejores data curriculum, architecture e inductive bias.
\end{knowledgebox}

\teachervoice{La clase no presentó BabyLM como «los niños demuestran que el scaling está equivocado», sino como una pregunta de investigación: si los humanos pueden aprender con menos datos lingüísticos, ¿podrían los modelos alcanzar una mayor sample efficiency mediante datos de mayor calidad, estructuras distintas o interacción?}

\subsection{Modelos pequeños y on-device model}

Los Minified LLM y los on-device model cambian el objetivo de «la máxima capacidad general» a «costo, latencia y privacidad con una capacidad aceptable». Los modelos más pequeños pueden ganar utilidad práctica mediante destilación, quantization, pruning, datos especializados y retrieval.

\lecturefigure{slide-056.jpg}{Motivaciones y enfoques de la miniaturización y de los LLM en el dispositivo}{slide deck oficial de CS25 V4, página 56}

\begin{knowledgebox}{Glosario: cuatro enfoques de compresión de modelos}
\begin{tabularx}{\textwidth}{L{0.20\textwidth} X}
\toprule
Método & Mecanismo central \\
\midrule
Distillation & Entrena un student más pequeño con las salidas, los logits o los reasoning traces del teacher. \\
Quantization & Representa weights/activations con un bit-width menor, lo que reduce memoria y cómputo. \\
Pruning & Elimina conexiones, channels, heads o blocks de baja contribución. \\
Specialization & Usa datos de dominio, adapters, tools o retrieval para acotar el rango de capacidades que deben quedar parametrizadas. \\
\bottomrule
\end{tabularx}
\end{knowledgebox}

\subsection{Memory augmentation y personalization}

Los parámetros de un LLM congelado no se actualizan automáticamente con cada conversación. La memory augmentation (aumento de memoria) coloca parte de la información en un store externo y la recupera hacia el context cuando hace falta; la personalization (personalización) permite que el sistema aproveche de forma continua las preferencias, el historial y las restricciones del usuario. Las dos diapositivas siguientes ponen lado a lado «modificar la entrada, modificar la memoria externa y modificar los parámetros»; al leerlas, primero hay que identificar qué objeto cambia cada enfoque.

\lecturefigure{slide-057.jpg}{Contradicción entre el conocimiento congelado y la necesidad de personalización a largo plazo}{slide deck oficial de CS25 V4, página 57}

\lecturefigure{slide-058.jpg}{Enfoques de adaptación: Fine-tuning, prompt, memory bank y RAG}{slide deck oficial de CS25 V4, página 58}

\begin{importantbox}{Los cuatro mecanismos de adaptación no son intercambiables}
\begin{tabularx}{\textwidth}{L{0.19\textwidth} L{0.24\textwidth} X}
\toprule
Mecanismo & Qué cambia & Problemas adecuados \\
\midrule
Prompt/context & La entrada de esta ocasión & Instrucciones de tareas temporales y poca evidencia. \\
Retrieval/RAG & Memory externa actualizable & Hechos, documentos, historial del usuario y evidencia citable. \\
Adapter/fine-tuning & Parte o la totalidad de los parámetros & Comportamiento estable, formato, distribución de dominio y adaptación de habilidades. \\
Model editing & Parámetros locales previamente localizados & Modificación dirigida de unos pocos hechos o asociaciones; requiere verificar efectos secundarios. \\
\bottomrule
\end{tabularx}
\end{importantbox}

\begin{importantbox}{Retrieval-augmented generation en su forma mínima}
Dada una query $q$, una colección de documentos $\mathcal{D}$, un recuperador $R$ y un generador $G$:
\[
z_{1:k}=R(q,\mathcal{D}),\qquad y=G(q,z_{1:k}).
\]
$z_{1:k}$ es la evidence top-$k$. La recuperación amplía el conocimiento accesible, pero no garantiza que la evidence sea correcta, que el orden sea razonable ni que el generator realmente la utilice.
\end{importantbox}

\begin{lstlisting}[caption={Flujo mínimo de recuperación desde memoria externa y generación}]
def answer(query, memory, encoder, generator, top_k=5):
    query_vector = encoder(query)
    evidence = memory.search(query_vector, top_k=top_k)
    prompt = build_prompt(query=query, evidence=evidence)
    return generator(prompt), evidence
\end{lstlisting}

\teachervoice{El ponente distinguió expresamente entre la memory externa y la skill interna: RAG sirve para complementar hechos e historial del usuario, pero no hace que el modelo adquiera automáticamente nuevas capacidades básicas. Si el datastore es de mala calidad, la recuperación solo llevará los errores más rápido al context.}

\subsection{Synthetic data y Phi: ¿puede la calidad de los datos sustituir parte de la escala?}

Los datos sintéticos (synthetic data) se generan con un modelo potente, que produce instructions, answers o reasoning traces, con los que después se entrenan otros modelos. Reducen el costo de la recolección manual, pero también pueden copiar los errores, los sesgos de estilo y los puntos ciegos del teacher. Esta subsección toma Phi como caso representativo del enfoque centrado en la calidad de los datos: primero revisa la generación y el filtrado de datos, y después el alcance de los resultados con modelos pequeños.

\lecturefigure{slide-059.jpg}{Datos sintéticos de preentrenamiento, filtrado y distillation}{slide deck oficial de CS25 V4, página 59}

\begin{warningbox}{Tres riesgos de ciclo cerrado de los synthetic data}
\begin{enumerate}
  \item \textbf{Error inheritance}: el student aprende de nuevo los errores del teacher;
  \item \textbf{Diversity collapse}: la distribución generada es más estrecha que el mundo real y, tras entrenamientos repetidos, los patrones convergen;
  \item \textbf{Evaluation leakage}: los prompts o las respuestas coinciden con el benchmark, lo que infla las puntuaciones.
\end{enumerate}
Por ello se necesitan provenance, deduplication, quality filter y comparación con datos reales.
\end{warningbox}

\lecturefigure{slide-060.jpg}{Resultados presentados en clase de Phi-2, que entrena un modelo pequeño con datos de alta calidad y sintéticos}{slide deck oficial de CS25 V4, página 60}

\teachervoice{El ponente situó el takeaway principal de Phi en la data source y la quality, en lugar de presentar el bajo número de parámetros como algo mágico. Que un modelo con menos parámetros conserve sus capacidades depende de la cobertura de tareas, del filtrado de datos y del alcance de la evaluación.}

\begin{knowledgebox}{Lectura de la figura: Phi representa el enfoque de que «cada token valga más»}
La clase destaca los textbook-quality data y los synthetic reasoning data, que permiten a un modelo más pequeño acercarse a modelos más grandes en algunas tareas. La figura no demuestra que los modelos pequeños sean equivalentes en todas las capacidades; muestra que, además de la escala de parámetros, la data selection y el curriculum pueden cambiar los resultados de manera significativa.
\end{knowledgebox}

\subsection{¿Learning o memorizing?}

Cuando los datos de entrenamiento y los benchmarks no son transparentes, la salida del modelo puede provenir de una generalización composicional o de una memorización aproximada. La contaminación de prueba (test contamination) hace que el modelo haya visto las preguntas o contenido muy similar antes de la evaluación, de modo que la puntuación deja de representar la generalización a tareas nuevas. La diapositiva del debate que sigue debe leerse como un problema de evaluation design, no como una exigencia de elegir con prisa entre «piensa» y «solo recita».

\lecturefigure{slide-061.jpg}{Conocimiento nuevo, combinación de patrones y el debate sobre la memorization en los LLM}{slide deck oficial de CS25 V4, página 61}

\begin{importantbox}{Tres niveles de una memorization audit}
\begin{enumerate}
  \item \textbf{Exact match}: si el modelo puede reproducir fragmentos largos del texto de entrenamiento;
  \item \textbf{Near duplicate}: si el conjunto de entrenamiento contiene ligeras reformulaciones de las preguntas, las plantillas o las respuestas;
  \item \textbf{Behavioral generalization}: si el mecanismo sigue funcionando al cambiar entidades, estructura o restricciones.
\end{enumerate}
Quedarse en el primer nivel deja pasar mucha contamination, y fijarse solo en la exactitud final no permite distinguir entre memorización y generalización.
\end{importantbox}

\begin{warningbox}{No confundir la controversia de derechos de autor con la controversia cognitiva}
Si el modelo «entiende» es una cuestión cognitiva y de comportamiento; si los datos de entrenamiento cuentan con autorización y si la salida reproduce en exceso la obra original son cuestiones legales y de gobernanza. Ambas están relacionadas, pero una conclusión no puede sustituir a la otra.
\end{warningbox}

\subsection{Continual learning: no es otro nombre del reentrenamiento periódico}

El continual learning (aprendizaje continuo) exige que el modelo se actualice cuando llegan nuevas experiencias, conservando a la vez sus capacidades previas y controlando el olvido catastrófico. La clase toma como contraste el aprendizaje cotidiano de las personas y critica que la trace distillation o el fine-tuning periódico se presenten directamente como automejora permanente. Esta subsección define primero el stability--plasticity tradeoff y después explica por qué mejorar la puntuación en una sola tarea nueva no basta para demostrar aprendizaje continuo.

\lecturefigure{slide-062.jpg}{Aprendizaje continuo, adaptación en línea y olvido catastrófico}{slide deck oficial de CS25 V4, página 62}

\begin{importantbox}{Continual-learning objective}
Para una secuencia de tareas $\mathcal{T}_1,\ldots,\mathcal{T}_m$, el modelo actualizado debe reducir la pérdida en la tarea nueva y, al mismo tiempo, limitar la degradación en las tareas anteriores:
\[
\min_\theta\ \mathcal{L}_{\mathcal{T}_m}(\theta)+
\lambda\sum_{i<m}\Omega_i(\theta,\theta_{i}^{*}).
\]
$\Omega_i$ es la regularization, el replay o la parameter constraint que preserva el conocimiento previo, y $\lambda$ controla el equilibrio entre plasticity y stability. Lo verdaderamente difícil es verificar el olvido cuando no se dispone de todos los datos anteriores.
\end{importantbox}

\teachervoice{El ponente ilustró el problema con un contraste intuitivo: las personas no se sientan cada dos meses a «volver a leer todo internet» para seguir aprendiendo. Destilar nuevos traces en el modelo se parece más a un reentrenamiento y no equivale a un aprendizaje continuo en línea y a lo largo de toda la vida.}

\subsection{Interpretabilidad y model editing}

La interpretabilidad busca entender cómo las representaciones y los cálculos internos producen la salida; la mechanistic interpretability se centra más en components, circuits y causal intervention concretos. El model editing, por su parte, pretende cambiar de forma dirigida una asociación factual concreta sin reentrenar todo el modelo.

\lecturefigure{slide-063.jpg}{La caja negra de los modelos grandes, el control y la mechanistic interpretability}{slide deck oficial de CS25 V4, página 63}

\begin{knowledgebox}{No mezclar tres tipos de «explicación»}
\begin{tabularx}{\textwidth}{L{0.24\textwidth} X}
\toprule
Nivel & Pregunta que responde \\
\midrule
Behavioral explanation & ¿Cómo influyen los cambios en la entrada sobre la salida y en qué distribuciones falla el modelo? \\
Attribution/probing & ¿Qué tokens, activations o features se relacionan con la salida? \\
Mechanistic/causal analysis & ¿Intervenir un component cambia de forma estable el comportamiento objetivo? \\
\bottomrule
\end{tabularx}
La correlación ayuda a localizar, pero no constituye por sí misma un mecanismo causal.
\end{knowledgebox}

\lecturefigure{slide-064.jpg}{Esquema de clase de ROME y del model editing factual}{slide deck oficial de CS25 V4, página 64}

\begin{importantbox}{Forma abstracta de un rank-one edit}
Si se quiere actualizar de forma dirigida la matriz de pesos $W$ de cierta capa, puede escribirse
\[
W' = W + uv^\top,
\]
donde $u,v$ son vectores obtenidos a partir de la key/value relation objetivo. Un rank-one update tiene bajo costo, pero su validación debe revisar a la vez edit success, paraphrase generalization, locality y unrelated knowledge damage.
\end{importantbox}

\begin{warningbox}{Corregir un hecho no significa que la «base de conocimiento» del modelo quede reparada}
Un hecho puede estar distribuido en varias capas y contextos; una edición local también puede alterar entidades cercanas o patrones lingüísticos. El model editing requiere una regression suite, no solo mostrar que el prompt objetivo ya produce la nueva respuesta.
\end{warningbox}

\subsection{MoE, modularity y autorreflexión}

El curso vuelve a MoE y plantea una pregunta más audaz: ¿podría un solo modelo formar módulos reutilizables en su interior, como las áreas funcionales del cerebro humano? Esta analogía sirve para estimular el diseño de architecture, pero no puede ocultar las dificultades reales del router, la expert specialization y la global coordination.

\lecturefigure{slide-065.jpg}{Estructura de routed experts de MoE y activación dispersa}{slide deck oficial de CS25 V4, página 65}

\lecturefigure{slide-066.jpg}{Propuesta abierta: de múltiples experts a la modularización dentro de un solo modelo}{slide deck oficial de CS25 V4, página 66}

\begin{warningbox}{La modularización necesita interfaces e indicadores de validación}
Si entre los módulos no hay entradas y salidas definidas, routing criterion, capacity y failure isolation, la «modularización» es solo un diagrama conceptual. Entre los indicadores medibles están la expert load, el token drop, la routing entropy, la specialization, el costo de comunicación y la propagación de fallas.
\end{warningbox}

La self-reflection (autorreflexión) hace que el modelo lea su propia respuesta, genere una critique y la modifique de forma iterativa. Puede aumentar el test-time compute y también sacar a la luz errores; pero si el critique model comparte puntos ciegos con el modelo original, el ciclo puede producir solo errores más largos y más seguros de sí mismos.

\lecturefigure{slide-067.jpg}{La self-reflection itera la salida mediante critique y revision}{slide deck oficial de CS25 V4, página 67}

\lecturefigure{slide-068.jpg}{Cambios de desempeño y límites de la self-reflection de varios niveles}{slide deck oficial de CS25 V4, página 68}

\begin{importantbox}{Reflection loop}
Sea la respuesta inicial $y_0=G(x)$; en la ronda $t$ se genera la critique $c_t=C(x,y_t)$ y después se revisa
\[
y_{t+1}=R(x,y_t,c_t).
\]
$G$ es el generador, $C$ es el proceso de critique y $R$ es el proceso de revision. El ciclo necesita una stopping rule externa y un verifier; de lo contrario, «seguir reflexionando» puede aumentar el costo sin aumentar la exactitud.
\end{importantbox}

\subsection{Hallucination, calibration y evidencia externa}

La diapositiva de hallucination describe el problema como «el modelo no sabe que no sabe». Dicho con más precisión, el modelo puede seguir produciendo lenguaje de alta confianza aunque carezca de evidencia. La calibration (calibración) se ocupa de si la confianza coincide con la tasa real de aciertos; retrieval y verifier intentan, por su parte, incorporar evidencia externa.

\lecturefigure{slide-069.jpg}{Hallucination, verificación de hechos, calibration y el enfoque de retrieval}{slide deck oficial de CS25 V4, página 69}

\begin{importantbox}{Definición de calibration}
Si el modelo asigna una confianza $q$ a un conjunto de respuestas, la calibración ideal exige
\[
P(\text{correct}\mid \text{confidence}=q)\approx q.
\]
Estar bien calibrado no equivale a tener alta exactitud; un modelo que responde «no estoy seguro» con frecuencia puede estar muy bien calibrado. A la inversa, un modelo de alta exactitud puede mostrar exceso de confianza en sus pocos errores.
\end{importantbox}

\begin{warningbox}{La temperature no es una perilla de factualidad}
Reducir la sampling temperature hace que la distribución sea más aguda y más cercana a la greedy decoding, pero no puede convertir un conocimiento erróneo en uno correcto. La factualidad requiere el apoyo conjunto de evidence, retrieval, verification, datos de entrenamiento y restricciones de la tarea.
\end{warningbox}

\subsection{Resumen de la sección}

La principal carencia de los LLM no es un bug aislado, sino un conjunto de problemas acoplados entre sí: los datos y el cómputo determinan el alcance de lo que puede aprenderse, la memory externa ofrece un canal de actualización, el aprendizaje continuo y el model editing modifican los parámetros, la interpretabilidad ayuda a localizar mecanismos, y la reflection y la calibration cambian el comportamiento en el momento de la inferencia. Ningún método por sí solo resuelve automáticamente a la vez el conocimiento, las habilidades, la factualidad y la adaptación a largo plazo.
\section{Reasoning scaffold: hacer que el proceso intermedio pueda buscarse y revisarse}

La última parte de contenido sobre modelos del curso trata Chain-of-Thought, Tree of Thoughts y Socratic Questioning. La idea común es cambiar «una salida directa» por un proceso intermedio de varios pasos, de modo que el modelo obtenga más test-time compute y el sistema conserve puntos donde revisar y retroceder.

\subsection{Por qué CoT puede ser útil}

Chain-of-Thought (cadena de pensamiento, CoT) pide al modelo, mediante ejemplos o instrucciones, que genere pasos intermedios. En tareas aritméticas y simbólicas de varios pasos es difícil asignar directamente la respuesta a la pregunta; al descomponer la tarea, la relación local de cada paso puede quedar más cerca de la distribución de entrenamiento. Las dos páginas siguientes muestran primero la intermediate guidance y después comparan la diferencia de información entre un standard prompt y un CoT example.

\lecturefigure{slide-070.jpg}{La intermediate guidance despliega una sola predicción en un razonamiento de varios pasos}{Slide deck oficial de CS25 V4, página 70}

\lecturefigure{slide-071.jpg}{Comparación entre prompting estándar y Chain-of-Thought prompting}{Slide deck oficial de CS25 V4, página 71}

\begin{importantbox}{Descomposición probabilística de CoT}
Si la reasoning trace intermedia es $r$ y la respuesta final es $y$, el proceso de generación puede escribirse como
\[
p(y\mid x)=\sum_r p(r\mid x)p(y\mid x,r).
\]
En la práctica, el decoding solo muestrea o busca unos cuantos $r$. La ventaja de CoT proviene de un cómputo intermedio y una descomposición más ricos; no garantiza que la trace muestreada refleje fielmente el proceso causal interno del modelo.
\end{importantbox}

\teachervoice{El ponente usa la multiplicación de números de diez cifras para explicar la intuición de los pasos intermedios: las personas también suelen calcular por pasos. El valor adicional de CoT es que los errores quedan visibles: el lector puede localizar en qué paso el modelo omitió o malinterpretó algo, en lugar de ver solo una respuesta incorrecta.}

\subsection{Más allá de la mejora de desempeño: hay que ver el tipo de error}

CoT suele producir mejoras en modelos grandes, pero todavía genera pasos irrelevantes, errores locales o conclusiones erróneas tras un proceso correcto. La clase divide los errores en missing step y semantic misunderstanding, y destaca en particular que los modelos pequeños pueden ser inestables incluso en el mapping local. En la lectura de las figuras de esta sección, lo central es la error taxonomy y el model size, no solo recordar que el desempeño promedio sube.

\lecturefigure{slide-072.jpg}{Resumen de la clase sobre la mejora de CoT y los errores restantes}{Slide deck oficial de CS25 V4, página 72}

\begin{warningbox}{Un rationale largo no equivale a un reasoning confiable}
El modelo puede generar explicaciones con formato completo pero sin relación con la respuesta, y también puede redactar la justificación hacia atrás después de obtener la respuesta. La validación debe examinar si los estados intermedios son ejecutables y verificables y si admiten un counterfactual change, no solo si el texto parece un razonamiento.
\end{warningbox}

\lecturefigure{slide-073.jpg}{Errores de missing-step y de semantic-understanding en modelos pequeños}{Slide deck oficial de CS25 V4, página 73}

\begin{knowledgebox}{Por qué los modelos pequeños fallan con más facilidad en CoT}
CoT aumenta la longitud de la secuencia y la coherencia que debe mantenerse. Si el modelo carece de capacidades semánticas o aritméticas básicas, más tokens crean más puntos donde equivocarse. Las rutas de mejora incluyen supervision por pasos, tool execution, verifier, distillation y una state representation más estructurada.
\end{knowledgebox}

\subsection{Generalized CoT: el razonamiento no tiene una sola línea}

El curso plantea que la definición de CoT es demasiado rígida: distintos problemas pueden requerir búsqueda con ramificaciones, retroceso, recursión sobre subproblemas o herramientas externas. Lo central del generalized reasoning es elegir el computation graph adecuado. Las tres páginas siguientes pasan de una trace lineal a la tree search y la recursive decomposition, y comparan state, transition y verifier.

\lecturefigure{slide-074.jpg}{De un Chain-of-Thought fijo hacia estructuras de razonamiento más generales}{Slide deck oficial de CS25 V4, página 74}

\begin{importantbox}{Tres ejes de diseño de un reasoning scaffold}
\begin{enumerate}
  \item \textbf{State}: si el estado intermedio es lenguaje natural, un programa, una ecuación o una observation del entorno;
  \item \textbf{Transition}: si el siguiente paso proviene de un muestreo del modelo, de reglas, de una búsqueda o de la ejecución de una herramienta;
  \item \textbf{Verifier}: cómo se puntúa, poda, retrocede y termina cada estado local.
\end{enumerate}
Solo cuando los tres quedan definidos con claridad, un método de razonamiento deja de ser un simple estilo de prompt.
\end{importantbox}

\subsection{Tree of Thoughts y Socratic decomposition}

Tree of Thoughts (árbol de pensamientos, ToT) mantiene varios estados candidatos y realiza lookahead/backtracking; Socratic Questioning (cuestionamiento socrático) plantea subproblemas de forma recursiva y luego combina sus respuestas en el problema original. Ambos intercambian más inference compute por espacio de búsqueda.

\lecturefigure{slide-075.jpg}{Ramificación, evaluación y retroceso en Tree of Thoughts}{Slide deck oficial de CS25 V4, página 75}

\begin{importantbox}{Búsqueda de anchura limitada}
Si cada estado se expande en $b$ candidatos, se conservan los top-$k$ y la profundidad es $d$, el número de llamadas al modelo crece aproximadamente como $O(dkb)$. El sesgo del verifier determina qué ramas se eliminan prematuramente; por ello, la calidad de la búsqueda no depende solo del generator, sino también del state scoring.
\end{importantbox}

\lecturefigure{slide-076.jpg}{Socratic Questioning genera y responde subproblemas de forma recursiva}{Slide deck oficial de CS25 V4, página 76}

\begin{knowledgebox}{Lectura de la figura: diferencia entre decomposition y search}
La decomposition divide un problema grande en problemas pequeños relativamente independientes; la search elige entre varias rutas candidatas. El Socratic method pone más énfasis en la estructura de subproblemas; ToT, en los estados candidatos y el retroceso. Los sistemas complejos suelen combinar ambos y llamar en los nodos hoja a un calculator, un code executor o a retrieval.
\end{knowledgebox}

\subsection{Resumen de la sección}

El valor de un reasoning scaffold está en aumentar el cómputo intermedio, exponer los errores y ofrecer interfaces de búsqueda y verificación. CoT, ToT y Socratic decomposition no producen automáticamente un reasoning genuino; solo se acercan a un proceso de cómputo controlable cuando se combinan con estados estructurados, tools, verifier y stopping rule.
\section{De Language Model a AI Agent: por qué una sola llamada no basta}

En el minuto 49, la clase pasa al tema de los agent. Esta transición no es un cambio de buzzword, sino un cambio en la unidad del sistema: un language model recibe tokens y genera tokens; un agent, además, tiene que mantener un state, leer la memory, elegir una action, observar el resultado en el entorno y seguir ejecutando según la retroalimentación. Una llamada al modelo puede ser un component del agent, pero no es un agent completo.

\subsection{Agent stack: qué más se necesita fuera del model}

La diapositiva de apertura enumera action, memory, communication y emergent architecture. La diapositiva Building AI Agents divide el problema en why y how: por qué se necesitan agent y cómo hacer que el modelo entre en interfaces reales.

\lecturefigure{slide-078.jpg}{Alcance del curso y líneas de investigación de la parte de AI Agent}{Diapositivas oficiales de CS25 V4, página 78}

\lecturefigure{slide-079.jpg}{Building AI Agents: marco del problema en términos de why y how}{Diapositivas oficiales de CS25 V4, página 79}

\begin{importantbox}{El ciclo mínimo de un agent}
Sea $s_t$ el estado del entorno; el modelo elige una action a partir de la observation $o_t$, la memory $m_t$ y el objetivo $g$:
\[
a_t=\pi_\theta(o_t,m_t,g),\qquad
s_{t+1}=\mathcal{E}(s_t,a_t),\qquad
m_{t+1}=U(m_t,o_t,a_t,s_{t+1}).
\]
$\pi_\theta$ es la policy, $\mathcal{E}$ es la transición del entorno y $U$ es la memory update. El agent tiene que procesar los resultados de sus acciones, no solo predecir el siguiente token.
\end{importantbox}

\lecturefigure{slide-080.jpg}{Una sola foundation-model call carece de estado a largo plazo y de retroalimentación de las acciones}{Diapositivas oficiales de CS25 V4, página 80}

\lecturefigure{slide-081.jpg}{La agent architecture agrega memory, tools y environment loop fuera del modelo}{Diapositivas oficiales de CS25 V4, página 81}

\begin{knowledgebox}{Términos clave: los seis componentes del agent stack}
\begin{tabularx}{\textwidth}{L{0.19\textwidth} X}
\toprule
Componente & Función \\
\midrule
Model/policy & Interpreta la observation y propone un plan o una action. \\
State & Guarda la tarea actual, las páginas, los archivos, las variables y el punto de ejecución. \\
Memory & Guarda información entre pasos o entre sesiones y permite recuperarla. \\
Tool/action layer & Asigna cada action estructurada a una API, un browser, un shell o un dispositivo. \\
Environment feedback & Devuelve éxito, errores, cambios en la página, resultados de pruebas u observation de sensores. \\
Verifier/guardrail & Verifica restricciones, permisos, resultados y condiciones de terminación. \\
\bottomrule
\end{tabularx}
\end{knowledgebox}

\teachervoice{El ponente explica de forma directa por qué se pasa a los agent: una sola foundation-model call no basta, porque las tareas reales también requieren memory, long context, personalization, actions, internet access y retroalimentación entre pasos.}

\subsection{Demos de Flight y driving-test: muestran action, no prueban confiabilidad}

Dos diapositivas de demo muestran a un agent temprano de MultiOn reservando un vuelo y completando el California online driving test. Muestran que el modelo puede convertir un objetivo en lenguaje natural en operaciones del navegador, pero un éxito único en un video público no permite inferir la tasa de éxito general, la seguridad de los pagos ni la robustez entre sitios web. Al leer la figura, primero identifica la observation, la action y la condición de término; después distingue la evidencia de demostración de la evidencia de despliegue.

\lecturefigure{slide-082.jpg}{Caso de reserva de vuelos con un AI agent mostrado en clase}{Diapositivas oficiales de CS25 V4, página 82}

\lecturefigure{slide-083.jpg}{Caso de agent para el online driving-test mostrado en clase}{Diapositivas oficiales de CS25 V4, página 83}

\begin{warningbox}{Cómo leer correctamente la demo evidence}
Una demo puede probar que «esta action trajectory ocurrió al menos una vez», pero no prueba la tasa de éxito promedio, la worst-case safety ni el despliegue sin supervisión. Para pasar de una capability demo a una engineering claim hacen falta un conjunto de tareas, el número de repeticiones, una clasificación de fallas, el alcance de los permisos, una estrategia de recuperación y una reproducción independiente.
\end{warningbox}

\teachervoice{El ponente llama a la demo del driving-test una de las primeras demostraciones de un real-world agent. Las notas conservan este hecho histórico, pero no lo presentan como la conclusión de que funcione de forma estable en cualquier sitio web, cuenta o dispositivo.}

\subsection{Por qué buscar una human-like interface}

El atractivo de un human-like agent es reutilizar las interfaces existentes: las páginas web, las aplicaciones de desktop y las mobile app ya están diseñadas para personas. Si el agent puede ver, hacer clic y escribir, no necesita esperar a que cada servicio ofrezca una API; además, puede aprender las preferencias del usuario a partir de sus demostraciones. Esa misma ventaja amplía el riesgo, porque el inicio de sesión, los pagos y los datos personales también quedan expuestos en el action space.

\lecturefigure{slide-084.jpg}{Un human-like agent reutiliza las interfaces, las preferencias y las demostraciones humanas}{Diapositivas oficiales de CS25 V4, página 84}

\begin{importantbox}{Interface coverage y safety boundary son dos caras de la misma moneda}
Un UI control más general aumenta la cobertura, pero debilita el schema, los permisos y los error code que proporciona una typed API. En el diseño conviene dar prioridad a las API con restricciones; solo cuando no exista una interfaz se amplía a UI control, y a las action de alto riesgo se les agregan confirmation, least privilege, audit log y rollback/recovery.
\end{importantbox}

\subsection{El autonomy level es una gradación de responsabilidad, no una categoría de marketing}

La clase toma prestada la clasificación L0--L5 de la conducción autónoma para hablar de agent: en los niveles bajos el control es humano y en los altos se retira gradualmente el human fallback. La analogía ayuda a explicar cómo cambia la responsabilidad, pero no sirve para clasificar directamente un automóvil o un producto de agent concreto. La siguiente figura debe leerse según «quién supervisa, quién puede tomar el control y quién responde cuando algo falla», no como una puntuación única de capacidad.

\lecturefigure{slide-085.jpg}{Analogía de cinco niveles, de human control a full autonomy}{Diapositivas oficiales de CS25 V4, página 85}

\begin{warningbox}{Los ejemplos de productos en clase son analogías informales}
El ponente asoció de palabra varios sistemas vehiculares con L4/L5. Las definiciones formales de SAE, el dominio de diseño operativo y el estado actual de los productos son mucho más rigurosos que la analogía de clase y cambian con el tiempo. La única conclusión didáctica que se conserva aquí es esta: retirar el human fallback eleva considerablemente los requisitos de verificación y de responsabilidad.
\end{warningbox}

\begin{knowledgebox}{Clasificación operativa de la agent autonomy}
\begin{tabularx}{\textwidth}{L{0.17\textwidth} X}
\toprule
Nivel & Una definición más adecuada para agent de software \\
\midrule
Assist & Ofrece sugerencias; no ejecuta action externas. \\
Approve-each-step & Cada action requiere la confirmación del usuario. \\
Bounded delegation & Ejecuta de forma automática dentro de un scope, un presupuesto y unas herramientas predefinidos. \\
Supervised autonomy & Puede ejecutar durante periodos largos, pero con monitor, interrupt y fallback en tiempo real. \\
Unsupervised autonomy & Completa tareas de alto impacto sin supervisión humana en línea; requiere los límites formales más estrictos y verificación independiente. \\
\bottomrule
\end{tabularx}
\end{knowledgebox}

\subsection{API y direct computer interaction}

Un agent tiene dos vías principales de action. La API ofrece parámetros estructurados, valores de retorno y permisos; la direct interaction opera las interfaces humanas mediante screen, keyboard, mouse o accessibility tree. La segunda tiene mayor cobertura, pero es más vulnerable al layout, la latency, los popup y la ambigüedad visual.

\lecturefigure{slide-087.jpg}{Dos vías: API control y direct computer interaction}{Diapositivas oficiales de CS25 V4, página 87}

\begin{importantbox}{Action schema}
Una action con capacidad de auditoría no debería ser solo texto libre, sino algo cercano a
\[
a_t=(\texttt{tool},\texttt{arguments},\texttt{precondition},\texttt{risk},\texttt{idempotency}).
\]
\texttt{precondition} describe el estado previo a la ejecución, \texttt{risk} determina si se requiere confirmación e \texttt{idempotency} indica si es seguro repetir la llamada. Las action estructuradas dan al verifier y a la recovery un objeto definido sobre el cual trabajar.
\end{importantbox}

\lecturefigure{slide-088.jpg}{Universal Action API y ejemplo de operación entre aplicaciones}{Diapositivas oficiales de CS25 V4, página 88}

\begin{warningbox}{«No se necesita API» no equivale a «la API no tiene valor»}
Un UI agent puede cubrir sistemas que no tienen API, pero una API estable suele ser más rápida, más barata y más fácil de autorizar y de probar. La arquitectura ideal es el capability routing: dar prioridad a las herramientas con restricciones y, cuando sea indispensable usar la UI, reducir el área visible, el action set y el credential scope.
\end{warningbox}

\subsection{Resumen de la sección}

El agent coloca al modelo de lenguaje dentro de un observation--action--feedback loop y le agrega state, memory, tool, verifier y recovery. Las demos de la clase mostraron la action capability, mientras que la autonomy y el UI control dejan al descubierto problemas de responsabilidad y de permisos. El verdadero límite del sistema no está en si «el modelo sabe hacer clic», sino en si las fallas pueden detectarse, acotarse y recuperarse.
\section{Neural Compute Unit, memoria de largo plazo y personalización}

Ver el LLM como una neural compute unit es la analogía de sistemas más importante de la segunda mitad de la clase. El modelo, como una unidad de cómputo, recibe token instructions y produce tokens, pero no obtiene de forma automática persistent memory, process isolation, device driver ni permission model. Solo al entender el alcance de la analogía se puede diseñar correctamente una agent architecture.

\subsection{Model as compute unit}

En la clase se comparan los tokens de lenguaje natural con las binary instructions. Lo que tienen en común es que la entrada pasa por un cómputo y se convierte en salida; la diferencia es que el comportamiento de un LLM es probabilístico, su espacio semántico es abierto y su entrada y salida carecen del tipado estricto y del determinismo de una ISA tradicional. Al leer la figura conviene mantener a la vez la compute analogy y las missing systems guarantees.

\lecturefigure{slide-090.jpg}{El AI model como neural compute unit token-in/token-out}{Presentación oficial de CS25 V4, página 90}

\begin{knowledgebox}{Lectura de la figura: límites de validez de la CPU analogy}
\begin{tabularx}{\textwidth}{L{0.22\textwidth} L{0.32\textwidth} >{\raggedright\arraybackslash}X}
\toprule
Nivel & Dónde la analogía es válida & Dónde la analogía falla \\
\midrule
Instruction & El prompt/context especifica el cómputo de esta ronda & El lenguaje natural es ambiguo y carece de opcodes fijos y de sistema de tipos. \\
Compute & Los parámetros ejecutan operaciones matriciales a gran escala & La salida es una distribución de probabilidad; no se pueden suponer garantías de determinismo ni de corrección. \\
Memory & El context ofrece un espacio de trabajo & El context tiene un límite de longitud y no se vuelve persistente de forma automática. \\
I/O & Las tools amplían las capacidades externas & Los permisos, las excepciones y los efectos secundarios deben gestionarse explícitamente en el sistema. \\
\bottomrule
\end{tabularx}
\end{knowledgebox}

\teachervoice{El ponente llama al modelo neural compute unit e imagina escribir un nuevo programming language para el Transformer. El valor didáctico de esta analogía está en convertir el prompt, la memory, las tools y las variables en componentes de un sistema, en lugar de atribuirle al modelo rasgos humanos como si fuera un cerebro completo.}

\subsection{Looped Transformer y cómputo repetido}

Un forward pass corresponde a un cómputo de profundidad fija. La diapositiva del Looped Transformer propone bloques compartidos o repetidos para que el modelo ejecute más iteraciones con los mismos parámetros. Se parece al reasoning loop: ambos intercambian más computation steps por un procesamiento más complejo, pero deben resolver la termination y la stability.

\lecturefigure{slide-091.jpg}{Looped Transformer y módulos de cómputo repetido}{Presentación oficial de CS25 V4, página 91}

\begin{importantbox}{Iterative computation}
Si el mismo bloque $F_\theta$ se repite $K$ veces:
\[
h^{(k+1)}=F_\theta(h^{(k)},x),\qquad k=0,\ldots,K-1.
\]
Compartir parámetros reduce la velocidad de crecimiento del modelo, pero puede traer problemas de optimization, de fixed-point stability y de stopping dinámico. Repetir más capas no equivale automáticamente a un reasoning más profundo.
\end{importantbox}

\subsection{La long-term memory no es solo vector search}

En la clase se compara la persistent memory con un disco: primero se codifican los documentos o las experiencias como embeddings y se almacenan, y cuando hace falta se recuperan hacia el context. Los sistemas actuales suelen usar nearest-neighbor search, pero una memoria de largo plazo real también requiere tiempo, jerarquía, relaciones y estrategias de actualización. Esta sección explica primero el embedding pipeline actual y después amplía el memory schema con las open questions planteadas en la clase.

\lecturefigure{slide-092.jpg}{Problemas de almacenamiento, embedding y recuperación en la long-term memory}{Presentación oficial de CS25 V4, página 92}

\begin{importantbox}{Memory retrieval score}
Para un query embedding $q$ y un memory item $m_i$, la similitud más simple puede escribirse como
\[
s_i=\frac{q^\top m_i}{\lVert q\rVert_2\lVert m_i\rVert_2},\qquad
\mathcal{I}=\operatorname{TopK}_i(s_i).
\]
Esto solo aprovecha la similitud semántica. Un memory ranking real también debería incorporar recency, importance, user scope, causal link, task state y access control.
\end{importantbox}

\begin{knowledgebox}{Concepto previo: los cuatro niveles de la agent memory}
\begin{enumerate}
  \item \textbf{Working memory}: el context actual, el scratchpad y los tool results;
  \item \textbf{Episodic memory}: interactions y action trajectories pasadas;
  \item \textbf{Semantic memory}: hechos estables, documentos y conocimiento del usuario;
  \item \textbf{Procedural memory}: workflows, policies y skills reutilizables.
\end{enumerate}
Construir una sola base de datos vectorial normalmente no basta para cubrir a la vez estas cuatro semánticas.
\end{knowledgebox}

\teachervoice{El ponente señala que la memory actual suele quedarse en un k-NN simple. La coherencia temporal, la graph structure, la online adaptation y los datos en cambio constante no quedan resueltos de forma natural; «conectar una vector database» es solo el punto de partida.}

\subsection{Personalization: las preferencias cambian y los datos tienen límites}

La diapositiva de Personalization aspira a que el agent sea una extensión digital del usuario. Para ello debe aprender su estilo de escritura, su agenda, su tolerancia al riesgo y sus objetivos de largo plazo, y al mismo tiempo distinguir entre preferencias estables, instrucciones de una sola vez e información sensible. Las dos diapositivas siguientes deben leerse considerando a la vez la capability y el data boundary, con especial atención a cómo se recopilan, actualizan, revocan y sobrescriben las preferencias.

\lecturefigure{slide-093.jpg}{Objetivos de la user-agent alignment y la personalization}{Presentación oficial de CS25 V4, página 93}

\lecturefigure{slide-094.jpg}{Retos de recopilación de preferencias, privacidad, cambio y alignment}{Presentación oficial de CS25 V4, página 94}

\teachervoice{El ponente recurre al hippocampus y a los datos personales en cambio constante para recordar al lector que la memory no es un repositorio estático de documentos. El sistema debe manejar el orden temporal, la estructura de relaciones, la actualización en línea y los cambios en las preferencias del usuario.}

\begin{warningbox}{Cuatro fallas comunes de la Personalization}
\begin{enumerate}
  \item Interpretar erróneamente una conducta aislada como una preferencia permanente;
  \item Mezclar la memory de distintos usuarios o workspaces;
  \item Guardar datos sensibles sin límite para «entender mejor al usuario»;
  \item Obedecer ciegamente cuando las preferencias del usuario entran en conflicto con la seguridad, la ley o la policy de la organización.
\end{enumerate}
El sistema necesita consent, scope, retention, edit/delete y policy precedence.
\end{warningbox}

\begin{importantbox}{La preference update requiere decaimiento temporal y confirmación}
El score de una preference $p$ puede escribirse como
\[
S(p)=\sum_j w_j\,e^{-\lambda\Delta t_j}\,\operatorname{evidence}_j(p),
\]
donde $w_j$ indica la confiabilidad de la evidencia, $\Delta t_j$ es el intervalo de tiempo y $\lambda$ controla el decaimiento. Las preferencias de alto riesgo no deben inferirse solo a partir de conductas implícitas; deben requerir confirmación explícita.
\end{importantbox}

\subsection{Resumen de la sección}

La neural-compute analogy aclara los límites entre model, context, memory y tool, pero no permite tratar al LLM como una CPU con una ISA determinista y garantías de aislamiento. La memoria de largo plazo necesita jerarquía, tiempo y permisos, y la personalización necesita mecanismos de consent, actualización y borrado. Cuanto más cerca del usuario están las capacidades del sistema, más se convierte la memory governance en un elemento central del diseño.
\section{Multi-Agent Systems: paralelizar es fácil; la comunicación y la corrección de errores son más difíciles}

La trayectoria larga de un solo agent puede ser lenta y difícilmente abarca tareas complejas. Un multi-agent system intenta repartir la tarea entre varios specialized worker y después consolidar los resultados mediante un manager. Lo más valioso de la clase no es «abrir más modelos», sino la secuencia de diagramas progresivos que muestra delegation, verification, conflict y redo.

\subsection{Por qué se necesitan varios agents}

La paralelización puede reducir el wall-clock time de las subtareas independientes, y la specialization permite que distintos worker usen herramientas o prompts diferentes. Esto supone que la tarea pueda descomponerse correctamente, que los resultados puedan combinarse y que el shared state no se sobrescriba de forma cruzada. Las dos diapositivas siguientes presentan primero la forma del sistema y después enumeran los beneficios del paralelismo y la specialization; solo entonces se aborda el costo del protocolo.

\lecturefigure{slide-096.jpg}{Varios autonomous agent forman un sistema colaborativo}{Diapositivas oficiales de CS25 V4, página 96}

\lecturefigure{slide-097.jpg}{Parallelization, specialization y la motivación de multi-agent}{Diapositivas oficiales de CS25 V4, página 97}

\begin{importantbox}{Límite superior de la ganancia por paralelismo}
Si la proporción paralelizable de la tarea es $p$, la aceleración ideal con $n$ worker está limitada por la Amdahl's law:
\[
\operatorname{Speedup}(n)=\frac{1}{(1-p)+p/n}.
\]
Los costos de descomposición, comunicación, verificación y combinación hacen que la aceleración real sea menor. El valor de un sistema multi-agent debe comprobarse con end-to-end latency, cost y success rate, no medirse por el número de agents.
\end{importantbox}

\teachervoice{El docente enfatiza que los beneficios directos de multi-agent son la parallelization y los specialized workers, pero de inmediato desplaza el foco hacia la communication: el lenguaje natural es ambiguo, y los agents también pueden malinterpretarse entre sí, igual que las personas.}

\subsection{Hierarchy y message protocol}

La clase propone una manager--worker hierarchy: el usuario conversa con el manager, el manager asigna tareas y mantiene el estado global, y los worker atienden objetivos locales. La jerarquía reduce las interfaces que enfrenta el usuario, pero concentra en el manager la descomposición correcta y la sincronización del estado. Las tres diapositivas siguientes van de la jerarquía abstracta al message y al state; al leer las figuras, conviene buscar primero la ownership y la version boundary.

\lecturefigure{slide-098.jpg}{Estructura jerárquica de comunicación entre manager y worker}{Diapositivas oficiales de CS25 V4, página 98}

\lecturefigure{slide-099.jpg}{Ambigüedad y problemas de sincronización en la natural-language communication}{Diapositivas oficiales de CS25 V4, página 99}

\begin{importantbox}{Un typed message es más fácil de verificar que el texto libre}
Una worker request debe incluir al menos
\[
m=(\texttt{task\_id},\texttt{goal},\texttt{inputs},\texttt{constraints},\texttt{deadline},\texttt{expected\_artifact}).
\]
El mensaje de respuesta debe incluir status, artifact, evidence, error y side effect. El lenguaje natural puede conservarse en el campo explanation, pero la información de control esencial debe estar estructurada.
\end{importantbox}

\lecturefigure{slide-100.jpg}{Diagrama inicial de comunicación entre el manager state y el worker state}{Diapositivas oficiales de CS25 V4, página 100}

\begin{knowledgebox}{Lectura de la figura: lo que debe sincronizarse es el state, no solo el historial de conversación}
El manager debe conocer la versión de la tarea, la instantánea de las entradas, las action que el worker ya ejecutó y la ubicación de los productos. Si dos worker trabajan sobre snapshots distintos, puede que no sea posible combinarlos aunque al final ambos textos «digan que terminaron». La versión del estado y la ownership importan más que la fluidez de la redacción.
\end{knowledgebox}

\subsection{Verification: una declaración de finalización no puede tomarse directamente como un hecho}

El cambio central del progressive diagram es la incorporación del verify step. Que un worker responda «done» es solo una claim; el manager o un verifier independiente debe revisar el artifact, las pruebas, el estado del entorno o evidencia externa. Los dos retained state siguientes muestran, respectivamente, una verification normal y un conflicto con la claim; lo importante es cómo el canal de evidencia modifica el flujo de control.

\lecturefigure{slide-103.jpg}{El manager aplica verification a la declaración de finalización del worker}{Diapositivas oficiales de CS25 V4, página 103}

\begin{importantbox}{Verification contract}
Para el producto $a$ de un worker y la especificación $c$, el verifier produce
\[
v=V(a,c,e)\in\{\texttt{pass},\texttt{fail},\texttt{unknown}\},
\]
donde $e$ es la evidencia de pruebas, registros, archivos, páginas o del entorno. \texttt{unknown} es importante: cuando la evidencia es insuficiente, no debe forzarse a pass.
\end{importantbox}

\lecturefigure{slide-104.jpg}{Scenario 1: la claim del worker entra en conflicto con el resultado de la verificación}{Diapositivas oficiales de CS25 V4, página 104}

\begin{warningbox}{El conflict handling no debe reducirse a «preguntarle otra vez al mismo modelo»}
Si el verifier usa solo el mismo contexto, el mismo modelo y la misma suposición errónea, la revisión puede repetir el error original. Una verificación más sólida debe cambiar el canal de evidencia: ejecutar pruebas, leer el estado real, usar otra tool, revisar el checksum o exigir una proof ejecutable.
\end{warningbox}

\subsection{Correction y redo path}

El diagrama de estado final incorpora el re-do: el manager envía al worker la causa del fallo y el context actualizado, y el worker vuelve a ejecutar la tarea. Una recovery eficaz debe distinguir entre transient failure, bad plan, permission error y side effects irreversibles.

\lecturefigure{slide-107.jpg}{Scenario 2: tras un fallo de verificación, se devuelve el context y se vuelve a ejecutar}{Diapositivas oficiales de CS25 V4, página 107}

\begin{lstlisting}[caption={Protocolo estructurado de verification y redo entre manager y worker}]
def run_task(manager, worker, verifier, task, max_attempts=3):
    context = manager.snapshot(task)
    for attempt in range(max_attempts):
        result = worker.execute(task, context=context)
        verdict = verifier.check(task.spec, result.artifact)
        if verdict.status == "pass":
            return manager.commit(result)
        context = manager.revise_context(
            previous=context,
            evidence=verdict.evidence,
            error=verdict.reason,
        )
    raise TaskFailed(task.id, context)
\end{lstlisting}

\begin{warningbox}{Antes del redo, determina si la action puede repetirse}
Enviar correos, hacer pagos, borrar archivos y enviar formularios pueden no ser operaciones idempotent. La recovery requiere un action ID, un side-effect log, una deduplication key y un compensation plan; de lo contrario, «volver a ejecutar» convierte un solo error en side effects duplicados.
\end{warningbox}

\teachervoice{En clase se usa el progressive diagram para explicar que el manager no puede confiar en la declaración de finalización del worker y debe verificarla; si el resultado es incorrecto, debe devolver el failure context y repetir la tarea. Este ciclo cerrado se acerca más a un sistema real que el lema de la «colaboración automática entre múltiples agents».}

\subsection{Resumen de la sección}

Los beneficios de un multi-agent system provienen del paralelismo y la specialization; los riesgos principales provienen de la descomposición, la comunicación, el estado compartido, la verificación y los side effects. Una manager--worker hierarchy solo resulta más confiable que un único agent cuando se acompaña de un typed protocol, evidence-based verification, versioned state y bounded redo.
\section{Confiabilidad, Plan Divergence y LLM OS}

Al final, el curso concentra los problemas pendientes del agent en la reliability. El software tradicional puede ejecutarse repetidamente con la misma entrada; la agent policy, en cambio, interactúa con un entorno dinámico a lo largo de muchos stochastic step. Aunque la tasa de éxito local parezca alta, en las trayectorias largas los fallos se acumulan con rapidez.

\subsection{Problemas centrales de los agents autónomos}

La diapositiva Key Issues enumera reliability, looping, plan divergence, benchmark, observabilidad, security y human fallback. No son funciones complementarias que se añaden después del despliegue, sino parte de la agent architecture.

\lecturefigure{slide-109.jpg}{Problemas de confiabilidad, ciclos, pruebas y seguridad de un autonomous agent}{Slide deck oficial de CS25 V4, página 109}

\begin{importantbox}{Efecto multiplicativo en la tasa de éxito de trayectorias largas}
Si cada paso tiene una probabilidad de éxito independiente $p$, la probabilidad de que una trajectory de longitud $T$ tenga éxito completo es aproximadamente
\[
P(\text{all success})=p^T.
\]
Incluso con $p=0.99$, al ejecutar 100 pasos la probabilidad de terminar sin ningún error es de apenas $0.99^{100}\approx 36.6\%$. Los errores reales no son independientes, y una desviación temprana además modifica las observation posteriores, lo que agrava el fallo.
\end{importantbox}

\teachervoice{El ponente usa la frase «aun con un agent al 95\%, el 5\% restante va a fallar» para explicar la diferencia en el despliegue de software: en un programa tradicional se puede corregir un deterministic bug, mientras que los pocos fallos de un stochastic agent reaparecen una y otra vez a escala de muchos usuarios y en trayectorias largas.}

\begin{knowledgebox}{Qué debe registrar la observabilidad}
\begin{enumerate}
  \item La observation, la selected action y la tool response de cada paso;
  \item Las versiones de model/prompt/tool y el memory snapshot;
  \item La permission decision, la user confirmation y el side effect;
  \item El verifier verdict, los retry, el fallback y el final outcome.
\end{enumerate}
El propósito de los registros es permitir la reproducción y la localización de fallos, no conservar contenido sensible de forma ilimitada; deben acompañarse de minimización, cifrado y una retention policy.
\end{knowledgebox}

\subsection{Plan divergence: cómo volver al rumbo después de desviarse}

La diapositiva de plan divergence contrasta el ideal path con el actual path. Si en algún paso el agent interpreta mal una página, una herramienta devuelve un error o la descomposición del objetivo es incorrecta, las action posteriores pueden apoyarse en un state equivocado y terminar en un loop. Esta sección retoma la reliability de la diapositiva anterior y explica con más detalle por qué el sistema debe comparar periódicamente el state previsto en el plan con la observation real.

\lecturefigure{slide-110.jpg}{Plan divergence: la trayectoria real se aleja paso a paso de la ruta ideal}{Slide deck oficial de CS25 V4, página 110}

\begin{importantbox}{Recovery policy}
El agent no debe preguntarse solo «cuál es el siguiente paso», sino también verificar periódicamente
\[
d_t=D(s_t,\hat{s}_t,g),
\]
donde $s_t$ es la observation real, $\hat{s}_t$ es el state esperado según el plan y $g$ es el objetivo. Cuando el divergence score $d_t$ supera un umbral, el sistema debe pausar, volver a observar, revertir a un checkpoint, replanificar o solicitar intervención humana.
\end{importantbox}

\begin{warningbox}{Un loop detector no puede limitarse a contar texto repetido}
El agent puede repetir la misma action fallida con otras palabras, o bien quedar en un ciclo entre páginas. La detección debe combinar action signature, state hash, error pattern, progress metric y consumo de presupuesto. Al alcanzar el límite de pasos, tiempo o costo, es obligatorio detenerlo de forma definitiva.
\end{warningbox}

\teachervoice{El ponente califica a AutoGPT como un buen prototype que, sin embargo, «en la práctica no puede hacer gran cosa», porque tiende a entrar en ciclos, desviarse y cometer errores de manera continua. Lo importante aquí no es descalificar el prototipo, sino señalar que un sistema de trayectorias largas sin error correction no puede recuperarse por sí solo con solo seguir generando.}

\subsection{LLM OS: una analogía de sistema útil pero incompleta}

El diagrama LLM OS de Karpathy trata al LLM como la CPU, a la context window como la RAM, al embedding store como el file system, a calculator, Python y terminal como software tradicional o ALU, y a las imágenes, el audio, el navegador y otros modelos como peripheral.

\lecturefigure{slide-111.jpg}{Correspondencia de componentes del LLM OS de Karpathy}{Slide deck oficial de CS25 V4, página 111}

\begin{knowledgebox}{Lectura de la figura: tabla de correspondencias del LLM OS}
\begin{tabularx}{\textwidth}{L{0.22\textwidth} L{0.24\textwidth} X}
\toprule
Sistema tradicional & Analogía en el LLM system & Garantías que aún faltan \\
\midrule
CPU & Foundation model & Ejecución determinista, instrucciones tipadas y corrección. \\
RAM & Context window & Aislamiento, consistencia de paginación y escritura controlada. \\
File system & Retrieval/memory store & Transacciones, permisos, versiones y semántica de borrado. \\
ALU/software & Calculator, Python, tools & Sandbox, límites de recursos y semántica de errores. \\
Peripheral/network & Vision, audio, browser, other agents & Autenticación, fronteras de confianza y validación de entradas. \\
\bottomrule
\end{tabularx}
\end{knowledgebox}

\begin{warningbox}{La palabra «OS» no aporta por sí sola las capacidades de un sistema operativo}
Un OS real se encarga de scheduling, memory protection, process isolation, permission, interrupt y recovery. Si un LLM orchestration framework carece de estos mecanismos, solo tiene la forma de un OS, sin sus garantías de seguridad y confiabilidad.
\end{warningbox}

\subsection{Generalized AI system y action engine}

La siguiente diapositiva conecta la chat interface del usuario, el action engine, los agents, la memory, las tools y el mundo exterior en un generalized system. El action engine se encarga de route, plan, permission y aggregation, y es el control plane más importante fuera del modelo.

\lecturefigure{slide-112.jpg}{Diagrama del sistema con chat interface, action engine, agents, tools y memory}{Slide deck oficial de CS25 V4, página 112}

\teachervoice{La visión final sitúa al action engine, detrás de la chat interface, como el centro de enrutamiento: reparte las tareas entre distintos agents y luego integra los resultados. Esa posición también debe asumir permisos, presupuesto, verificación y cancelación, y no puede ser solo un prompt router.}

\begin{importantbox}{Control plane y data plane}
\begin{itemize}
  \item \textbf{Control plane}: interpreta el objetivo, selecciona agent/tool, asigna permisos, mantiene el state y decide retry/fallback;
  \item \textbf{Data plane}: ejecuta efectivamente las action de search, code, browser, database, message o device.
\end{itemize}
Al separar ambos, se puede impedir que una model proposal se convierta directamente en un side effect real.
\end{importantbox}

\subsection{Future needs: error correction, permission y sandbox}

La última diapositiva no concluye que «los agents ya están resueltos», sino que enumera la infraestructura pendiente: error correction, better framework, security, user permission, risky-domain stability y sandboxing. Al leer la figura, conviene tratar estos puntos como deployment precondition y relacionarlos uno por uno con la long-horizon failure, los permisos y los side effect vistos antes.

\lecturefigure{slide-113.jpg}{Los AI agents aún requieren error correction, security, permission y sandbox}{Slide deck oficial de CS25 V4, página 113}

\begin{importantbox}{Pila mínima de seguridad para action de alto riesgo}
\begin{enumerate}
  \item \textbf{Least privilege}: otorgar solo la credential y el action scope mínimos necesarios para la tarea actual;
  \item \textbf{Confirmation gate}: los pagos, envíos, eliminaciones, publicaciones y cambios de permisos requieren confirmación explícita;
  \item \textbf{Sandbox}: el código, los archivos y el acceso a la red se ejecutan en un entorno aislado;
  \item \textbf{Budget}: limitar tokens, pasos, tiempo, costo y llamadas externas;
  \item \textbf{Audit and recovery}: registrar los efectos secundarios y permitir la cancelación, la compensación o la toma de control humana.
\end{enumerate}
\end{importantbox}

\teachervoice{La conclusión final del curso no es que «los agents pronto serán completamente automáticos», sino que la error correction, la security, la user permission y el sandbox siguen sin resolverse. Sobre todo en escenarios de alto riesgo como finance o legal, la estabilidad y la seguridad deben anteponerse a una autonomy más amplia.}

\subsection{Resumen de la sección}

La agent reliability está dominada por el efecto multiplicativo de las trayectorias largas y por la plan divergence. El LLM OS ofrece un lenguaje para organizar los componentes, pero no produce por sí solo isolation, transaction ni recovery. Un sistema desplegable necesita que control plane, least privilege, verification, budget, fallback y sandbox restrinjan en conjunto la model proposal.
\section{Síntesis y ampliación}

Esta clase de Overview avanza desde el mecanismo de selección de información de la attention hasta los límites de seguridad de los agent. En apariencia abarca muchos temas, pero en el fondo puede reducirse a tres expansiones de frontera: primero, el Transformer permite que cada token lea directamente el contexto relevante; segundo, el LLM amplía el conjunto de tareas que puede abordar mediante escala, datos y post-training; tercero, el agent añade fuera del modelo memory, tool, state, communication y environment feedback. Cada expansión aporta nuevas capacidades y, al mismo tiempo, desplaza la verificación desde una sola salida hacia el sistema completo.

\subsection{Diez conclusiones centrales}

\begin{importantbox}{Del mecanismo del modelo a la ingeniería de sistemas}
\begin{enumerate}
  \item La attention selecciona el value mediante query--key matching, y la multi-head proporciona varios subespacios de representación;
  \item El Transformer cambia un emparejamiento global de $O(n^2)$ por rutas de información cortas y paralelismo en el entrenamiento;
  \item El LLM sigue basándose en el next-token objective; la instruction following y la safety provienen de entrenamiento adicional y de restricciones del sistema;
  \item Una emergence curve puede reflejar la scale, pero también puede verse amplificada por un metric threshold;
  \item RLHF, DPO, MoE y data quality son diseños de distintos niveles y no pueden reducirse a «un modelo más grande»;
  \item La memory externa, RAG, continual learning y model editing modifican objetos distintos;
  \item CoT, ToT y Socratic decomposition añaden cómputo intermedio, pero deben acompañarse de un verifier;
  \item Un agent es un observation--action--feedback loop, no una sola llamada al modelo;
  \item La dificultad principal de un sistema multi-agent está en el state, el protocol, la verification y la recovery, no en el arranque en paralelo;
  \item El LLM OS es un mapa de componentes; un despliegue real requiere además permission, isolation, budget, fallback y sandbox.
\end{enumerate}
\end{importantbox}

\subsection{Una tabla que articula toda la clase}

\begin{center}
\begin{longtable}{L{0.17\textwidth} L{0.22\textwidth} L{0.24\textwidth} L{0.25\textwidth}}
\toprule
Nivel & Capacidad nueva & Modo de falla nuevo & Verificación necesaria \\
\midrule
\endhead
Embedding & Representación continua compartida & La similitud no equivale a una relación factual & Pruebas de vecinos cercanos, sesgo y distribución \\
Attention & Lectura del contexto según su contenido & Pesos difíciles de interpretar directamente; costo cuadrático en secuencias largas & ablation, causal test, pruebas de eficiencia \\
LLM scale & Mayor cobertura de tareas y few-shot behavior & Costo, contaminación, capacidades impredecibles & benchmark multidimensional, auditoría de datos \\
Preference tuning & Mayor apego a las preferencias humanas u organizacionales & reward hacking, sesgo de preferencias & held-out preference, evaluación de seguridad \\
External memory & Hechos actualizables y personalización & Recuperación errónea, acceso sin autorización, memory caducada & Verificación de la recuperación, la evidencia y los permisos \\
Reasoning scaffold & Más test-time compute y pasos visibles & Coherente pero erróneo, costo de búsqueda & Herramientas, verificación y pruebas contrafactuales \\
Agent loop & Ejecución de tareas externas de largo plazo & side effect, loop, plan divergence & Bitácora de estado, presupuesto, reversión y recuperación \\
Multi-agent & Paralelismo y specialization & Ambigüedad de protocolo, conflictos, efectos secundarios duplicados & Mensajes estructurados, control de versiones y revisión independiente \\
\bottomrule
\end{longtable}
\end{center}

\subsection{Las preguntas que realmente deja el cierre de la clase}

La closing discussion del ponente concentra el futuro de los agent en error correction, security, permission y sandbox. Quedan además varias preguntas abiertas que vale la pena seguir:

\begin{enumerate}
  \item ¿Cómo medir el long-horizon success, en lugar de medir solo un tool call de un paso?
  \item ¿Cómo lograr personalization sin almacenar indefinidamente datos privados?
  \item ¿Cómo hacer que la memory admita a la vez tiempo, jerarquía, relaciones y olvido?
  \item ¿Cómo distinguir el valor explicativo de una reasoning trace del mecanismo causal real?
  \item ¿Cómo lograr que la multi-agent verification sea independiente del canal que originó el error?
  \item ¿Cómo definir estrategias de ejecución distintas para las action revocables, compensables e irreversibles?
  \item ¿Cómo hacer que los modelos más pequeños, los modelos especializados y la retrieval colaboren en el equilibrio entre capacidad y costo?
\end{enumerate}

\begin{warningbox}{Juicio de ingeniería final}
La action proposal que entrega el modelo debe separarse del side effect real. Todo sistema que involucre dinero, identidad, publicación, eliminación, permisos u operaciones irreversibles debe colocar el modelo detrás de un control plane restringido; una autonomy sin confirmation, audit, recovery ni sandbox no debe presentarse como automatización confiable.
\end{warningbox}

\subsection{Preguntas de autoevaluación}

\begin{enumerate}
  \item ¿Por qué el scaling por $\sqrt{d_k}$ de la attention ayuda a estabilizar el softmax?
  \item ¿En qué se diferencian las fuentes de Q/K/V en la self-attention y en la cross-attention?
  \item ¿Por qué una emergent metric puede presentar un salto aunque la capacidad subyacente varíe de forma suave?
  \item ¿En qué paso utilizan RLHF y DPO, respectivamente, la preference data?
  \item ¿Por qué difieren el número total de parámetros de un MoE y sus active parameters?
  \item ¿Qué modifica cada uno de RAG, fine-tuning, continual learning y model editing?
  \item ¿Por qué los pasos intermedios legibles de CoT no son necesariamente una explicación fiel?
  \item ¿Qué agent component le faltan a una sola LLM call?
  \item ¿Qué riesgos adicionales tiene el UI control frente al API control?
  \item ¿Por qué el manager no puede confiar directamente en el ``done'' del worker?
  \item Cuando una action no puede repetirse, ¿qué mecanismo debe añadirse al redo loop?
  \item En la analogía del LLM OS, ¿qué garantías de un OS tradicional siguen faltando?
\end{enumerate}

\subsection{Lecturas adicionales}

\begin{itemize}
  \item \href{https://arxiv.org/abs/1706.03762}{Vaswani et al., Attention Is All You Need}: el artículo original del Transformer y la multi-head attention.
  \item \href{https://arxiv.org/abs/2206.07682}{Wei et al., Emergent Abilities of Large Language Models} y \href{https://arxiv.org/abs/2304.15004}{Schaeffer et al., Are Emergent Abilities of Large Language Models a Mirage?}: para entender, respectivamente, el fenómeno y la metric caveat.
  \item \href{https://arxiv.org/abs/2305.18290}{Rafailov et al., Direct Preference Optimization}: el DPO objective y su relación con RLHF.
  \item \href{https://arxiv.org/abs/2202.05262}{Meng et al., Locating and Editing Factual Associations in GPT}: ROME y la factual model editing.
  \item \href{https://arxiv.org/abs/2201.11903}{Wei et al., Chain-of-Thought Prompting} y \href{https://arxiv.org/abs/2305.10601}{Yao et al., Tree of Thoughts}: la rationale lineal y el reasoning basado en búsqueda.
  \item \href{https://docs.google.com/presentation/d/1oXPs3LXtIVIsVbwTyGjAWj_aWvak9c1uNC4uhkS6glk/edit}{CS25 V4 official slide deck}: la canonical source de 114 diapositivas de esta clase.
\end{itemize}

\end{document}
